%% Auto-generated by md2tex.py — Springer Nature sn-jnl template
\documentclass[pdflatex,sn-mathphys-num]{sn-jnl}

\usepackage{graphicx}
\usepackage{multirow}
\usepackage{amsmath,amssymb,amsfonts}
\usepackage{amsthm}
\usepackage{mathrsfs}
\usepackage[title]{appendix}
\usepackage{xcolor}
\usepackage{textcomp}
\usepackage{manyfoot}
\usepackage{booktabs}
\usepackage{algorithm}
\usepackage{algorithmicx}
\usepackage{algpseudocode}
\usepackage{listings}
\usepackage{tikz}
\usetikzlibrary{shapes.geometric,arrows.meta,positioning,fit,calc,backgrounds}
\usepackage[utf8]{inputenc}
\usepackage{microtype}
\usepackage[T5]{fontenc}

\theoremstyle{thmstyleone}
\newtheorem{theorem}{Theorem}
\newtheorem{proposition}[theorem]{Proposition}
\theoremstyle{thmstyletwo}
\newtheorem{example}{Example}
\newtheorem{remark}{Remark}
\theoremstyle{thmstylethree}
\newtheorem{definition}{Definition}

\raggedbottom

\begin{document}
\sloppy

\title{Energy-Efficient Underwater Object Detection via Hierarchical Federated Learning and Parameter-Efficient Tuning}

% Tác giả 1: Nguyen Hai Nam (HUST)
\author[1]{\fnm{Hai Nam} \sur{Nguyen}}\email{nam.nghai24@gmail.com}

% Tác giả 2: Nguyen Thi My Binh (Corresponding Author - Hanoi University of Industry)
\author*[2]{\fnm{Thi My Binh} \sur{Nguyen}}\email{binhdungminhkhue@gmail.com}

% Tác giả 3: Pham Hai Van (HUST)
\author[1]{\fnm{Hai Van} \sur{Pham}}\email{haipv@soict.hust.edu.vn}

% Tác giả 4: Luong Ho Viet Duc (HUST) - Đã đảo lại theo yêu cầu
\author[1]{\fnm{Viet Duc Luong} \sur{Ho}}\email{luong.hovietduc@hust.edu.vn}

% Affiliations
\affil[1]{\orgdiv{School of Information and Communication Technology}, \orgname{Hanoi University of Science and Technology}, \orgaddress{\city{Hanoi}, \country{Vietnam}}}

\affil[2]{\orgname{Hanoi University of Industry}, \orgaddress{\country{Vietnam}}}

\abstract{The deployment of the Internet of Underwater Things (IoUT) is pivotal for ocean exploration. Integrating deep-learning-based object detection at the edge enhances the autonomy of autonomous underwater vehicles (AUVs). However, applying traditional Federated Learning (FL) to complex 2D vision tasks in IoUT faces severe physical constraints: ultra-limited acoustic bandwidth and strict battery budgets. Furthermore, standard gradient sparsification disrupts the spatial correlation of images, causing accuracy collapse, while directly averaging Low-Rank Adaptation (LoRA) matrices creates subspace misalignment. To overcome these limitations, this paper proposes FedKDL, an energy-optimized three-tier Hierarchical Federated Learning (HFL) architecture tailored for underwater object detection. FedKDL introduces a Dual-End Split Computing structure. At the seabed, AUVs train a compact student model using LoRA and Delta INT8 quantization, bounding the uplink payload to a minimal constant. At the surface, the gateway executes a novel LoRA-Projection Knowledge Distillation (KD) by extracting low-rank projections from an Teacher to synchronize knowledge without causing out-of-memory errors. To preserve spatial optimization trajectories, intermediate relay nodes perform SVD-LoRA Aggregation, mathematically eliminating cross-term errors. Finally, a nearest feasible association and HFL-Nearest cooperative routing mitigate depth-habitat non-IID data challenges. Extensive experiments on the URPC dataset demonstrate that FedKDL preserves state-of-the-art detection accuracy, maintains 100\% swarm connection coverage, and significantly reduces acoustic communication energy, offering a highly practical paradigm for large-scale underwater artificial intelligence.}

\keywords{Internet of Underwater Things, Hierarchical Federated Learning, Object Detection, Knowledge Distillation, Low-Rank Adaptation, Acoustic Communication, Edge AI.}

\maketitle

\section{Introduction}\label{sec:introduction}

Sự bùng nổ của mạng Internet vạn vật dưới nước (Internet of Underwater Things - IoUT) đã và đang tái định hình toàn diện các phương thức khảo sát đại dương, giám sát rạn san hô sinh thái và thực thi an ninh hàng hải. Trong hệ sinh thái phân tán này, các phương tiện lặn tự hành (Autonomous Underwater Vehicles - AUV) không còn đóng vai trò là những trạm thu thập dữ liệu thô sơ đơn lẻ, mà đang chuyển mình mạnh mẽ thành các thực thể cảm biến thông minh hoạt động theo bầy đàn. Nhằm nâng cao tính tự chủ tại thực địa, việc tích hợp các mô hình học sâu để thực hiện tác vụ nhận diện đối tượng 2D (Object Detection) ngay tại biên mạng là một nhu cầu cấp thiết, cho phép AUV phân loại các mục tiêu sinh học theo thời gian thực mà không cần đợi can thiệp từ trạm mặt nước.

Trong các chiến dịch khảo sát đại dương, mặc dù AUV dễ dàng thu thập khối lượng khổng lồ hình ảnh, việc gán nhãn hộp giới hạn (bounding-box) lại vô cùng đắt đỏ. Gần đây, Học tự giám sát (Self-Supervised Learning - SSL) \cite{kotthapalli2025self} nổi lên như một giải pháp đột phá, cho phép các mô hình thị giác tự học hỏi đặc trưng và sinh nhãn giả (pseudo-label) trực tiếp từ dữ liệu thô. Tuy nhiên, để triển khai quy trình này trên diện rộng, Học liên kết (Federated Learning - FL) \cite{mcmahan2017communication, zhang2021survey} là mảnh ghép viễn thông bắt buộc nhằm đồng bộ tri thức qua sóng âm mà không truyền tải hình ảnh nhạy cảm \cite{victor2022federated}. Xa hơn nữa, sự cộng hưởng giữa năng lực tự học của SSL và cơ cấu phân phối của hệ thống FL được đề xuất sẽ thiết lập nền tảng vững chắc cho một đường ống (pipeline) huấn luyện khép kín, tự chủ học hỏi liên tục dưới đại dương mà hoàn toàn không cần sự can thiệp của con người.

Tuy nhiên, việc triển khai FL truyền thống vào môi trường IoUT vấp phải những rào cản vật lý vô cùng khắc nghiệt. Kênh truyền sóng âm dưới nước bị thống trị bởi băng thông cực kỳ hạn hẹp ($\approx 10\text{--}15$ kbps), độ trễ lan truyền tín hiệu lớn do vận tốc âm thanh chậm, và nguồn pin sinh tồn của các nút cảm biến là hữu hạn \cite{nguyen2021federated, dang2024survey}. Việc truyền tải trực tiếp các mô hình nhận diện tích chập hiện đại chứa hàng chục triệu tham số \cite{tian2026yolov12} lên mặt nước sẽ ngay lập tức gây tắc nghẽn viễn thông \cite{niknam2020federated}. Để giảm tải áp lực này, các kiến trúc Học liên kết phân cấp (Hierarchical FL - HFL) \cite{liu2020client} kết hợp cơ chế hợp tác chọn lọc tại các trạm trung chuyển đã được đề xuất \cite{omeke2026energy}. Mặc dù vậy, một khoảng trống công nghệ khổng lồ vẫn tồn tại: hệ thống FL dưới nước hiện thời mới chỉ khả thi trên các tác vụ phân tích chuỗi thời gian 1D siêu nhẹ. 

Khi nỗ lực nâng cấp hệ thống để gánh vác tác vụ thị giác máy tính 2D, các phương pháp nén truyền thống lập tức bộc lộ điểm yếu chí mạng. Cách tiếp cận sử dụng kỹ thuật nén gradient thưa thớt hóa Top-K (Sparsification) \cite{lin2017deep} thất bại thảm hại do việc loại bỏ ngẫu nhiên 95\% gradient đã phá hủy hoàn toàn tính tương quan không gian (spatial correlation) của bản đồ đặc trưng, khiến độ chính xác mAP sụp đổ. Hướng đi khả dĩ hơn là ứng dụng Thích nghi Hạng thấp (LoRA) \cite{hu2022lora} để đóng băng mạng gốc. Dẫu vậy, việc đưa LoRA vào FL lại làm nảy sinh mâu thuẫn toán học cốt lõi: thuật toán FedAvg lấy trung bình trực tiếp trên các ma trận $A$ và $B$ riêng lẻ gây ra sự sai lệch không gian con (subspace misalignment) do sự thiếu hụt các tích chéo \cite{wang2024flora, bai2024federated}. Thêm vào đó, khi cố gắng áp dụng Chưng cất Tri thức (KD) \cite{hinton2015distilling} để bù đắp giới hạn của mô hình nhỏ, các kỹ thuật KD bảo toàn tương quan không gian truyền thống yêu cầu so khớp ma trận Gram khổng lồ \cite{tung2019similarity}, trực tiếp gây quá tải bộ nhớ (Out-of-Memory) cho thiết bị nhúng AUV. Hiện tượng trôi dạt trọng số do dữ liệu phi đồng nhất (Non-IID) \cite{li2020federated} càng đẩy mô hình toàn cục đến ranh giới phân kỳ.

Để phá vỡ hoàn toàn các nút thắt vật lý và thuật toán này, bài báo đề xuất \textbf{FedKDL} — một khung Học liên kết phân cấp nhận thức tri thức, thiết lập cấu trúc xử lý bất đối xứng hai bờ (Dual-End Split Computing) chuyên biệt cho ảnh 2D dưới đại dương. Các đóng góp mang tính đột phá bao gồm:

\begin{itemize}
\item \textbf{Kiến trúc phân tán với LoRA-Projection KD và Delta INT8:} Dưới đáy biển, AUV đóng băng mạng gốc, chỉ huấn luyện LoRA và dùng lượng tử hóa vi phân \textbf{Delta INT8} để cố định dung lượng gói tin ở mức rất nhỏ. Tại mặt nước, Gateway sử dụng kỹ thuật \textbf{LoRA-Projection} để chưng cất tri thức qua các phép chiếu hạng thấp ($h=Ax$), giúp đồng bộ mô hình mà không làm quá tải bộ nhớ thiết bị.

\item \textbf{Khắc phục sai lệch không gian bằng SVD-LoRA Aggregation:} Thay vì dùng FedAvg thông thường, các trạm trung chuyển (Relay) khôi phục ma trận cục bộ ($W_i = B_i A_i$), cộng dồn theo trọng số dữ liệu, rồi dùng Phân tích giá trị suy biến (SVD) để tái phân rã thành ma trận LoRA toàn cục. Việc này giúp triệt tiêu lỗi chéo và bảo toàn quỹ đạo tối ưu hóa.

\item \textbf{Định tuyến tối ưu và Hợp tác láng giềng:} Để giải quyết dữ liệu không đồng nhất (Non-IID) theo độ sâu, AUV sẽ kết nối với trạm trung chuyển có chi phí năng lượng phát sóng thấp nhất. Đồng thời, các trạm trung chuyển liên tục chia sẻ và nội suy mô hình với \textbf{trạm láng giềng gần nhất} để lan tỏa tri thức liên vùng. Cơ chế này giúp triệt tiêu hiện tượng trôi dạt mô hình (Client Drift) mà không tốn pin cho các thuật toán phức tạp.
\end{itemize}

Phần còn lại của bài báo được tổ chức như sau: Chương \ref{sec:related_work} tóm tắt các công trình liên quan và khoảng trống nghiên cứu. Chương \ref{sec:system_model} mô tả mô hình hệ thống, kênh truyền sóng âm và bài toán năng lượng. Chương \ref{sec:methodology} trình bày chi tiết thuật toán FedKDL. Chương \ref{sec:experiments} thảo luận kết quả mô phỏng toàn diện. Cuối cùng, Chương \ref{sec:conclusion} đưa ra kết luận.

\section{Related Work}\label{sec:related_work}

Việc triển khai Học liên kết (FL) trong mạng Internet vạn vật dưới nước (IoUT) là một giao điểm học thuật phức tạp giữa khoa học máy tính phân tán và vật lý viễn thông âm học. Mặc dù các khía cạnh nền tảng của FL—từ giới hạn vô tuyến, quyền riêng tư, phân mảnh dữ liệu đến tối ưu hóa năng lượng—đã được hệ thống hóa toàn diện qua nhiều nghiên cứu tổng quan \cite{niknam2020federated, zhang2021survey, nguyen2021federated, victor2022federated, tan2022towards, dang2024survey, yurdem2024federated, yang2025federated}, việc tích hợp chúng vào môi trường đại dương vẫn vấp phải những giới hạn vật lý khắc nghiệt. Dựa trên các rào cản đặc thù này, các nỗ lực nghiên cứu hiện hành được phân mảnh thành bốn định hướng cốt lõi sau:

\paragraph{Cấu trúc Mạng lưới và Học liên kết Phân cấp (Topology \& Hierarchical FL)}
Kiến trúc mạng hình sao (Star Topology) truyền thống của FL bộc lộ những điểm yếu chí mạng về độ trễ và tắc nghẽn băng thông khi đối mặt với quy mô lớn. Để giải quyết nút thắt này, mô hình Học liên kết Phân cấp (HFL) ba tầng Client-Edge-Cloud đã được chứng minh là một cơ chế định tuyến dữ liệu ưu việt, cho phép giảm tải lưu lượng diện rộng thông qua các trạm tổng hợp trung gian \cite{liu2020client, wu2024topology, li2025optimization}. 

Đưa cấu trúc này vào môi trường IoUT, các nghiên cứu đã tập trung khắc phục tính di động của AUV thông qua đồ thị liên kết động \cite{he2025mobility} và tối ưu hóa đa mục tiêu độ trễ - năng lượng bằng phương pháp cập nhật bất đồng bộ \cite{meng2022efficient}. Đáng chú ý, cơ chế phân cụm học đa tác vụ \cite{shu2022clustered} và thuật toán tổng hợp nhận thức năng lượng \cite{tyagi2026energy} đã mở ra hướng đi mới trong việc quản lý tài nguyên. Đặc biệt, chiến lược hợp tác có chọn lọc (Selective Cooperative Aggregation) \cite{omeke2026energy} đã thành công trong việc duy trì vòng đời của mạng cảm biến dưới nước bằng cách cắt giảm khắt khe các phiên giao tiếp liên cụm dư thừa. 

Song song với đó, các cấu trúc phi tập trung (Decentralized FL) và cơ chế bất đồng bộ cũng được nghiên cứu nhằm duy trì liên kết mạng trong điều kiện không tồn tại máy chủ trung tâm cố định \cite{shu2022clustered, he2025mobility, wu2024topology}. Tuy nhiên, một điểm mù công nghệ của các khung HFL dưới nước hiện thời là sự giới hạn về chiều dữ liệu. Các hệ thống này mới chỉ được kiểm chứng tính khả thi trên dữ liệu chuỗi thời gian (1D) với các mô hình phát hiện bất thường mang cấu trúc siêu nhẹ \cite{omeke2026energy}. Khi nỗ lực nâng cấp lên các tác vụ thị giác không gian đa chiều, tải lượng tham số lập tức vượt qua ngưỡng Shannon của kênh truyền sóng âm, làm sụp đổ hoàn toàn các cơ cấu định tuyến nói trên.

\paragraph{Xử lý Dữ liệu Phi đồng nhất (Data Heterogeneity)}
Sự phân bố của các quần thể sinh vật biển phụ thuộc chặt chẽ vào độ sâu và điều kiện thủy văn, dẫn đến việc dữ liệu thu thập bởi từng AUV mang tính chất phi đồng nhất (Non-IID) cực đoan. Trong điều kiện này, thuật toán nền tảng FedAvg \cite{mcmahan2017communication} đối mặt với hiện tượng trôi dạt trọng số cục bộ (client drift), khiến mô hình toàn cục mất khả năng hội tụ.

Để kiềm chế sự phân kỳ này, các hướng tiếp cận can thiệp trực tiếp vào hàm mục tiêu cục bộ đã được phát triển mạnh mẽ. Điển hình, FedProx \cite{li2020federated} bổ sung một số hạng phạt cận kề (proximal term) để neo giữ trọng số; SCAFFOLD \cite{karimireddy2020scaffold} sử dụng các biến kiểm soát (control variates) nhằm triệt tiêu phương sai cập nhật; trong khi MOON \cite{li2021model} ứng dụng học đối chiếu ở cấp độ mô hình (model-contrastive learning) để kéo gần không gian biểu diễn cục bộ với mô hình toàn cục. Song song đó, xu hướng cá nhân hóa mô hình thông qua học đa tác vụ \cite{smith2017federated} hoặc giao tiếp bằng nguyên mẫu (prototypes) thay vì truyền tải toàn bộ tham số mạng như FedProto \cite{tan2022fedproto} cũng cho thấy sự ưu việt trong việc bảo tồn các đặc trưng cục bộ.

Ở quy mô định tuyến, cơ chế gom cụm thiết bị (Clustering) nổi lên như một giải pháp dung hòa. Các thuật toán áp dụng K-Means trên đặc trưng thống kê phân phối \cite{alsenani2024fedsikd}, hoặc tối ưu hóa khoảng cách Earth Mover’s Distance (EMD) \cite{shu2022clustered} cho phép nhóm các AUV có miền dữ liệu tương đồng để học đa tác vụ, giảm thiểu xung đột tri thức trong quá trình tổng hợp.

Dẫu đạt được nền tảng toán học vững chắc, một rào cản chí mạng xuất hiện khi áp dụng các giải pháp trên vào thiết bị cảm biến biên. Việc tích hợp các ràng buộc chính quy hóa (regularization constraints) hay học đối chiếu đa lớp buộc thiết bị phải thực thi các chu trình tính toán vi phân ngược (backpropagation) khổng lồ ngay tại phần cứng cục bộ. Đối với các AUV chạy bằng nguồn pin giới hạn, sự gia tăng đột biến của khối lượng tính toán (FLOPs) này sẽ vắt kiệt năng lượng sinh tồn trước khi mô hình kịp hội tụ, khiến chúng trở nên phi thực tế trong mạng IoUT.

\paragraph{Nén Truyền thông và Thích nghi Hạng thấp (Communication Compression \& PEFT)}
Để đối phó với dải băng thông sóng âm cực hẹp, việc nén dung lượng gói tin cập nhật là yêu cầu bắt buộc. Kỹ thuật thưa thớt hóa gradient dựa trên độ lớn (Top-K Sparsification) như DGC \cite{lin2017deep} hay cắt tỉa mạng thích ứng \cite{liu2024adaptive} thường được áp dụng để cắt giảm tới 99\% lượng tham số truyền đi. Tuy nhiên, cách tiếp cận này thất bại nghiêm trọng khi áp dụng cho các tác vụ thị giác 2D phân giải cao. Việc loại bỏ ngẫu nhiên các gradient theo ngưỡng vô hướng sẽ phá hủy hoàn toàn cấu trúc hình học và tính tương quan không gian (spatial correlation) của bản đồ đặc trưng, gây mất mát vĩnh viễn biên dạng vật thể. Thay vào đó, kỹ thuật Thích nghi Hạng thấp (LoRA) \cite{hu2022lora} mở ra hướng đi tối ưu hơn bằng cách đóng băng mạng gốc khổng lồ và chỉ huấn luyện một không gian con (subspace) mang kích thước siêu nhỏ.

Dẫu vậy, việc đưa LoRA vào môi trường FL phân tán vấp phải một mâu thuẫn toán học cốt lõi. Thuật toán FedAvg truyền thống tiến hành cộng trung bình trực tiếp trên các ma trận $A$ và $B$ riêng lẻ, sinh ra sai lệch không gian con (subspace misalignment) do sự thiếu hụt các tích chéo $\frac{1}{2}(B_1+B_2)\frac{1}{2}(A_1+A_2) \neq \frac{1}{2}(B_1A_1+B_2A_2)$ \cite{wang2024flora, bai2024federated}. Để bảo toàn quỹ đạo tối ưu, các kiến trúc SOTA mới nhất buộc phải thay đổi phương thức tổng hợp: áp dụng Phân tích giá trị suy biến (SVD) để tái phân rã ma trận hiệu dụng $W = BA$ \cite{wu2022communication, bai2024federated}, sử dụng phân rã trực giao QR \cite{zhou2025ilora}, hoặc đóng băng bán phần $A$ \cite{sun2024improving}. Hơn nữa, khi kết hợp SVD với các chiến lược lượng tử hóa nhận thức lỗi (Quantization-aware) \cite{hu2025fedqlora}, lượng tử hóa theo ma trận Hessian \cite{su2024haflq}, hoặc Ternary weights \cite{he2025feddt}, hệ thống có thể ép chặt dung lượng truyền tải mà không gây ra hiện tượng phân kỳ do lỗi làm tròn (quantization bias).

Bên cạnh nén trọng số, Chưng cất Tri thức (KD) \cite{hinton2015distilling} cung cấp một lộ trình song song bằng cách thiết lập giao tiếp thông qua ``nhãn mềm'' (soft targets) trên dữ liệu proxy \cite{li2019fedmd, lin2020ensemble}. Tuy nhiên, các kỹ thuật KD bảo toàn độ tương đồng (Similarity-Preserving) truyền thống \cite{tung2019similarity} yêu cầu tính toán ma trận Gram trên toàn bộ khối lượng feature maps, gây quá tải RAM (Out-of-Memory) nghiêm trọng cho các thiết bị nhúng. Để khả thi hóa việc chuyển giao tri thức qua mạng vô tuyến, định hướng SOTA hiện tại (như AdaLD \cite{zhang2025communication}) đã chứng minh sức mạnh của việc trích xuất và đối sánh các phép chiếu ẩn hạng thấp ($h=Ax$) trực tiếp từ adapter LoRA thay vì logit truyền thống, tối thiểu hóa lượng byte cần truyền tải trong khi vẫn bảo toàn được biểu diễn không gian sâu của mô hình.

\paragraph{Ứng dụng Thị giác Máy tính trong IoUT (Computer Vision \& IoUT)}
Sự tiến hóa của các kiến trúc phát hiện đối tượng thời gian thực đã đạt đến trạng thái SOTA (State-of-the-Art) với các mô hình lấy cơ chế chú ý làm trung tâm (Attention-Centric) như YOLOv12 \cite{tian2026yolov12}. Việc tích hợp các mạng nơ-ron khổng lồ này vào môi trường Học liên kết thường đòi hỏi những chiến lược chia tách module phức tạp (như giữ lại Head, chỉ truyền tải Backbone và Neck) \cite{li2025ultraflwr} hoặc hoán đổi chéo mạng giải mã (Decoder) để vượt rào cản khác biệt miền (domain gap) \cite{jiang2025federated}.

Tuy nhiên, khi đặt vào môi trường dưới nước, các giải pháp thị giác FL hiện tại mới chỉ dừng ở mức sơ khai do sự hạn chế nghiêm ngặt của kênh truyền. Để thích ứng với băng thông sóng âm, các hệ thống buộc phải sử dụng các kiến trúc mạng siêu nhẹ (như U-Net) kết hợp bù nhiễu quang học \cite{yuan2025fedexchange}, hoặc áp dụng cắt tỉa ngẫu nhiên và lượng tử hóa thô ngay trên AUV \cite{wang2024underwater}. Thậm chí, khi cố gắng triển khai chưng cất tri thức tại các thiết bị biên (Edge AI) giới hạn tài nguyên, các nghiên cứu vẫn phải đánh đổi bằng việc gọt bỏ toàn bộ backbone nguyên bản và thay thế bằng các phiên bản rút gọn như MobileNet \cite{setyanto2025knowledge}, làm suy giảm vĩnh viễn năng lực trích xuất đặc trưng sâu của hệ thống.

\paragraph{Khoảng trống Nghiên cứu và Sự Đột phá của FedKDL (Research Gap)}
Tổng hợp từ bức tranh học thuật trên, một khoảng trống công nghệ sâu sắc vẫn đang tồn tại: \textit{Chưa có một kiến trúc nào đủ khả năng đưa các mô hình nhận diện đối tượng vượt qua kênh truyền sóng âm đại dương mà không làm vỡ cấu trúc không gian ảnh 2D}. 

Các phương pháp nén Top-K \cite{lin2017deep} làm sụp đổ mAP do phá hủy đặc trưng hình học, trong khi các thuật toán phân cụm chống Non-IID \cite{shu2022clustered, alsenani2024fedsikd} hoặc KD cục bộ lại vắt kiệt quỹ năng lượng sinh tồn của AUV do đòi hỏi điện toán lan truyền ngược khổng lồ. Mặt khác, các nỗ lực thích nghi hạng thấp (LoRA) \cite{hu2022lora} lại vướng phải sai lệch toán học khi tổng hợp \cite{wang2024flora, bai2024federated}, khiến mạng mất tính đồng bộ.

Khung giải pháp \textbf{FedKDL} được đề xuất trong nghiên cứu này là mảnh ghép hoàn hảo để giải quyết triệt để đa mâu thuẫn đó. Bằng cách thiết lập cấu trúc bất đối xứng hai bờ: dời toàn bộ thuật toán Chưng cất Tri thức bằng phép chiếu LoRA (LoRA-Projection KD) \cite{zhang2025communication} lên trạm Gateway mặt nước, đồng thời áp dụng lượng tử hóa Delta INT8 và quy tắc tổng hợp SVD-LoRA (SVD-LoRA Aggregation) \cite{bai2024federated} tại các nút trung chuyển (Relay), FedKDL hiện thực hóa việc bảo tồn sức mạnh nhận diện của mạng nơ-ron sâu dưới đáy biển với chi phí viễn thông và hao tổn điện toán ở mức cực tiểu.


\section{System Model and Problem Formulation}\label{sec:system_model}

\subsection{System Model}\label{subsec:system_model}

Việc triển khai nền tảng Học liên kết phân tán trong môi trường đại dương chịu sự chi phối tuyệt đối của các định luật vật lý về truyền âm. Khác với sóng vô tuyến trên mặt đất, tín hiệu âm thanh dưới nước bị suy hao nghiêm trọng theo khoảng cách và tần số, kết hợp với các nguồn nhiễu sinh học và động lực học phức tạp. Những rào cản này trực tiếp giới hạn dung lượng truyền tải (tốc độ Shannon) và tiêu tốn lượng lớn năng lượng phát sóng của thiết bị \cite{victor2022federated, omeke2026energy}. 

Do đó, trước khi thiết lập cấu trúc liên kết mạng và phát biểu bài toán tối ưu tổng thể, hệ thống cần được định lượng chặt chẽ thông qua các mô hình thành phần. Các mô hình vật lý này đóng vai trò là cơ sở toán học, thiết lập giới hạn và điều kiện biên cho mọi luồng dữ liệu (model updates) di chuyển giữa hai điểm nút bất kỳ trong mạng IoUT.

\subsubsection{Mô hình truyền tin âm thanh dưới nước (Underwater Acoustic Channel Model)}\label{sec1.4.2}

Trái ngược với các mạng vô tuyến (RF) trên mặt đất vốn sở hữu băng thông rộng và môi trường lan truyền tương đối đồng nhất, kênh truyền sóng âm dưới nước (Underwater Acoustic - UWA) là một trong những môi trường viễn thông khắc nghiệt nhất. Sự truyền tải dữ liệu tại đây chịu sự chi phối tuyệt đối của hiện tượng suy hao hấp thụ đa đường, dải băng thông hẹp và phổ nhiễu nền động lực học \cite{nguyen2021federated, victor2022federated}. Trong kiến trúc FedKDL, mô hình vật lý kênh truyền này đóng vai trò là "lưới lọc" đầu tiên, thiết lập các ranh giới toán học nghiêm ngặt để xây dựng đồ thị khả thi $\mathcal{G}$ và định lượng chi phí năng lượng cho toàn bộ vòng lặp Học liên kết \cite{tyagi2026energy, omeke2026energy}.

Khi một tín hiệu âm thanh di chuyển xuyên qua cột nước đại dương, năng lượng của nó bị suy giảm do hai nguyên nhân cốt lõi: sự lan tỏa hình học (geometric spreading) và sự chuyển hóa năng lượng âm thành nhiệt năng (absorption). Đối với một liên kết giữa nút phát $u$ và nút thu $v$ cách nhau một khoảng cách Euclidean 3D $d_{uv}$ (tính bằng mét), mức suy hao truyền dẫn $TL$ ở tần số sóng mang $f$ (kHz) được mô hình hóa theo công thức kinh nghiệm Thorp \cite{meng2022efficient, wang2024underwater}:
\begin{equation}\label{eq:thorp_tl}
TL(d_{uv}, f) = 10 k \log_{10}(d_{uv}) + \alpha(f) \frac{d_{uv}}{1000} \quad \text{(dB)}
\end{equation}
Trong đó, $k$ là hệ số lan truyền hình học (thường $k = 1.5$ đại diện cho môi trường lan truyền hình trụ trong nước biển nông). Hệ số hấp thụ nhiệt $\alpha(f)$ là một hàm phi tuyến tính phụ thuộc trực tiếp vào tần số, được biểu diễn bởi:
\begin{equation}\label{eq:absorption}
\alpha(f) = \frac{0.11f^2}{1+f^2} + \frac{44f^2}{4100+f^2} + 2.75 \times 10^{-4}f^2 + 0.003 \quad \text{(dB/km)}
\end{equation}
Phương trình \eqref{eq:absorption} giải thích một giới hạn vật lý chí mạng: khi tần số $f$ tăng, hệ số hấp thụ $\alpha(f)$ tăng vọt \cite{wang2024underwater}. Điều này trực tiếp giới hạn tần số hoạt động của các modem âm thanh IoUT ở mức chỉ vài chục kHz, kéo theo rào cản vật lý bất khả kháng về băng thông khả dụng cực hẹp cho hệ thống.

Khác với nhiễu trắng tĩnh trong mạng vô tuyến, nhiễu nền đại dương $NL$ là sự chồng chất của nhiều hiện tượng thủy động lực học biến thiên theo tần số. Tổng mức nhiễu nền tích lũy trong băng thông khả dụng $B$ (Hz) của máy thu được mô hình hóa theo phân tích của Wenz \cite{wang2024underwater}:
\begin{equation}\label{eq:wenz_noise}
NL(f, B) = 10 \log_{10}\left(\sum_{x \in \{\theta,s,w,th\}} 10^{N_x(f)/10}\right) + 10 \log_{10}(B) \quad \text{(dB re 1 $\mu$Pa)}
\end{equation}
Trong đó, các thành phần nhiễu phổ mật độ $N_x(f)$ được chia thành bốn dải tần thống trị. Cụ thể, nhiễu nhiễu loạn ($N_{\theta}$) chi phối ở dải tần cực thấp dưới 10 Hz với $N_{\theta}(f) = 17 - 30\log_{10}(f)$. Nhiễu tàu bè ($N_s$) gây ra bởi hoạt động hàng hải chiếm ưu thế ở dải 10--100 Hz và phụ thuộc vào hệ số mật độ giao thông $s \in [0, 1]$. Đáng chú ý nhất, nhiễu gió sóng ($N_w$) là nguồn nhiễu chủ đạo trong dải 100 Hz đến 100 kHz, trùng khớp hoàn toàn với dải tần hoạt động của IoUT và bị chi phối bởi vận tốc gió $w$. Cuối cùng, nhiễu nhiệt ($N_{th}$) sẽ thống trị ở các tần số cao trên 100 kHz \cite{tyagi2026energy, wang2024underwater}.

Để máy thu $v$ có thể giải mã thành công chuỗi bit trọng số mô hình từ nút phát $u$, tín hiệu âm thanh tại điểm thu phải vượt qua phổ nhiễu nền để đạt được một Tỷ số Tín hiệu trên Nhiễu mục tiêu ($\gamma_{tgt}$). Dựa trên phương trình Sonar thụ động, Tỷ số SNR thực tế trên liên kết $(u, v)$ được xác định \cite{he2025mobility, wang2024underwater}:
\begin{equation}\label{eq:sonar_snr}
SNR_{uv} = SL_u - TL(d_{uv}, f) - NL(f, B) - IL
\end{equation}
Với $IL$ là suy hao do lắp đặt và phần cứng. Từ phương trình \eqref{eq:sonar_snr}, ta suy ra mức công suất bức xạ âm thanh tối thiểu ($SL_u^{min}$) mà nút phát $u$ buộc phải tạo ra để thiết lập liên kết thành công:
\begin{equation}\label{eq:min_source_level}
SL_u^{min}(u, v) = \gamma_{tgt} + TL(d_{uv}, f) + NL(f, B) + IL
\end{equation}

Do giới hạn vật lý của đầu dò âm thanh, công suất phát tối đa của mỗi AUV bị khóa chặt tại ngưỡng $SL_{max}$. Hệ thống không cho phép các AUV cố gắng kết nối vô vọng ngoài tầm phủ sóng. Do đó, một Đồ thị khả thi không gian tĩnh (Feasibility Graph) $\mathcal{G}$ được thiết lập, chỉ bao gồm các liên kết thỏa mãn điều kiện công suất \cite{he2025mobility}:
\begin{equation}\label{eq:feasibility_graph}
\mathcal{G} = \bigl\{(u, v) \;\big|\; SL_u^{min}(u, v) \le SL_{max}\bigr\}
\end{equation}
Ràng buộc này là nền tảng cốt lõi để FedKDL thiết kế các thuật toán định tuyến và hợp tác liên cụm (Relay Cooperation) tại Phần IV, đảm bảo mạng lưới không bao giờ đưa ra các quyết định cập nhật viễn vông gây rớt gói tin.

Chỉ đối với các liên kết $(u, v) \in \mathcal{G}$ thỏa mãn tính khả thi, tốc độ truyền tải thông tin hiệu dụng mới được thiết lập. Tuân theo định lý Shannon-Hartley, tốc độ kênh truyền tĩnh $R_{uv}$ được xác định bởi \cite{omeke2026energy, wang2024underwater}:
\begin{equation}\label{eq:shannon_rate}
R_{uv} = B \log_2\left(1 + 10^{\gamma_{tgt}/10}\right) \quad \text{(bps)}
\end{equation}
Giá trị $R_{uv}$ này chính là nút thắt cổ chai lớn nhất của mạng IoUT. Nó chi phối trực tiếp đến độ trễ truyền thông, từ đó quyết định phần lớn tổng lượng năng lượng tiêu thụ $E_{tx}$ mà mạch vô tuyến của AUV phải gánh chịu trong quá trình tham gia Học liên kết.
\subsubsection{Mô hình Tính toán và Độ trễ (Computation and Latency Model)}\label{subsubsec:comp_latency_model}

Trong môi trường IoUT, sự kết hợp giữa vận tốc lan truyền sóng âm $c_s \approx 1500$ m/s (chậm hơn hàng trăm nghìn lần so với sóng điện từ) và băng thông âm học siêu hẹp tạo ra những thách thức chí mạng về tính thời gian thực. Theo nguyên lý thực thi song song (parallel execution) của hệ thống phân tán, độ trễ toàn vòng (round-trip latency) không đơn thuần là phép cộng dồn tuyến tính, mà được quyết định bởi "đường găng" (critical path) hay nhánh truyền tải chậm nhất trong mạng lưới \cite{he2025mobility, tyagi2026energy}. Độ trễ tổng thể $\tau_{round}^{(t)}$ được cấu thành từ hai pha chính: pha điện toán tại thiết bị biên và pha truyền thông dọc theo kiến trúc phân cấp. Do sự chênh lệch năng lực phần cứng, quá trình tổng hợp tại trạm Relay và máy chủ Gateway diễn ra gần như tức thời và không tạo ra độ trễ đáng kể so với môi trường truyền âm. Ngược lại, tại tầng đáy (Deep Layer), mỗi phương tiện tự hành (AUV) $i \in \mathcal{N}$ được trang bị một máy tính nhúng với năng lực vi xử lý bị giới hạn nghiêm ngặt. Nhiệm vụ của nút cảm biến là thực thi quá trình học máy nội bộ (Local Training) trên tập dữ liệu hình ảnh thu thập được $\mathcal{D}_i$. Bằng cách ánh xạ khối lượng tính toán qua quy mô của mô hình mạng nơ-ron, tổng số FLOPs yêu cầu cho một vòng lặp huấn luyện cục bộ tại thiết bị $i$ được xác định theo công thức \eqref{eq:phi_local}:
\begin{equation}\label{eq:phi_local}
\Phi_i = n_i \cdot E_{local} \cdot \left( \mu \cdot \text{FLOPs}_{sample} \right)
\end{equation}
Trong phương trình trên, đại lượng $n_i$ là số lượng mẫu ảnh đại diện cho quy mô tập dữ liệu cục bộ thực tế mà thiết bị biên $i$ đang nắm giữ, còn $E_{local}$ là số kỷ nguyên huấn luyện nội bộ (local epochs) diễn ra trên từng thiết bị. Đồng thời, $\text{FLOPs}_{sample}$ đại diện cho số lượng phép toán dấu phẩy động cơ sở cần thiết để lan truyền xuôi một bức ảnh qua mạng học sâu, và $\mu$ là hệ số nhân thuật toán đặc trưng cho quá trình lan truyền ngược (backward pass multiplier). Việc tham số hóa trực tiếp qua số lượng FLOPs phản ánh chính xác gánh nặng điện toán thực tế của các mạng tích chập phức tạp thay vì chỉ ước lượng thô thông qua số lượng trọng số $|\Theta|$ của mô hình.

Vì các AUV đồng loạt khởi chạy quá trình học máy, hệ thống sử dụng thời gian tính toán lớn nhất làm hằng số gốc cho vòng lặp. Độ trễ huấn luyện cục bộ (Computation Latency), ký hiệu là $\tau_{comp}^{(t)}$, phụ thuộc trực tiếp vào khối lượng phép toán $\Phi_i$ và xung nhịp CPU tĩnh ($f_i^{fix}$) của thiết bị \cite{meng2022efficient}, được xác định bởi phương trình:
\begin{equation}\label{eq:lat_comp}
\tau_{comp}^{(t)} = \frac{\Phi_i}{f_i^{fix}} \quad \text{(giây)}
\end{equation}

Bước sang pha truyền thông kênh truyền (Communication Latency), thời gian để một gói tin di chuyển trên liên kết vật lý $(u, v)$ bất kỳ là sự tổng hợp của trễ truyền dẫn (thời gian dữ liệu được đẩy vào kênh) và trễ lan truyền vật lý trong môi trường nước \cite{wang2024underwater}:
\begin{equation}\label{eq:link_latency}
\tau_{comm}(S; u, v) = \underbrace{\frac{S}{R_{uv}}}_{\text{Transmission Delay}} + \underbrace{\frac{d_{uv}}{c_s}}_{\text{Propagation Delay}} \quad \text{(giây)}
\end{equation}
Trong đó, $S$ là kích thước gói tin được gửi (bit), $R_{uv}$ là tốc độ Shannon tĩnh (bps) của liên kết, và $d_{uv}$ là khoảng cách hình học giữa nút phát và thu. 

Độ trễ truyền thông toàn mạng được phân rã thành ba chặng (hop-by-hop) nối tiếp nhau. Chặng thứ nhất đánh dấu quá trình luân chuyển dữ liệu từ AUV lên trạm trung gian (Uplink AUV-to-Relay). Đối với mỗi trạm Relay $m$, do các AUV thành viên sử dụng cấp phát đa truy cập và truyền song song, trạm trung gian buộc phải chờ đợi thiết bị có thời gian truyền lâu nhất trong cụm. Độ trễ chặng này đối với nhánh $m$ là:
\begin{equation}\label{eq:lat_a2r}
\tau_{a2r, m}^{(t)} = \max_{i \in \mathcal{C}_m} \Big( \tau_{comm}(S; i, m) \Big)
\end{equation}
Chặng thứ hai xuất hiện khi hệ thống kích hoạt cơ chế hợp tác liên cụm \cite{omeke2026energy}. Mỗi trạm Relay $m$ sẽ tiêu tốn thêm một khoảng thời gian để luân chuyển và vay mượn tri thức từ một đối tác láng giềng $m'$, với độ trễ được tính bằng:
\begin{equation}\label{eq:lat_coop}
\tau_{r2r, m}^{(t)} = \mathbb{1}\Big(m \in \mathcal{M}_{coop}^{(t)}\Big) \cdot \tau_{comm}(S; m, m')
\end{equation}
Cuối cùng, ở chặng thứ ba (Uplink Relay-to-Gateway), tất cả các Relay đồng loạt đẩy mô hình đại diện lên mặt nước. Gateway chỉ có thể tiến hành tổng hợp toàn cục khi nhận được gói tin từ trạm Relay xa nhất. Nút cổ chai tại Gateway tạo ra một hằng số trễ chung cho toàn mạng ở chặng này:
\begin{equation}\label{eq:lat_r2g}
\tau_{r2g}^{max} = \max_{j \in \mathcal{M}} \Big( \tau_{comm}(S; j, \text{Gateway}) \Big)
\end{equation}

Tổng hợp lại, chuỗi độ trễ truyền thông mà một nhánh Relay $m$ phải gánh chịu là sự cộng dồn của $\tau_{a2r, m}^{(t)}$, $\tau_{r2r, m}^{(t)}$ và $\tau_{r2g}^{max}$. Bầy đàn AUV chỉ hoàn tất một vòng lặp khi nhánh "ì ạch" nhất hoàn thành toàn bộ chu trình truyền tải. Do đó, tổng độ trễ của một vòng lặp Học liên kết $\tau_{round}^{(t)}$ được xác định chính xác bằng cách kết hợp pha điện toán nội bộ với đường găng truyền thông:
\begin{equation}\label{eq:total_latency}
\tau_{round}^{(t)} = \tau_{comp}^{(t)} + \max_{m \in \mathcal{M}} \Big( \tau_{a2r, m}^{(t)} + \tau_{r2r, m}^{(t)} \Big) + \tau_{r2g}^{max}
\end{equation}
Phương trình \eqref{eq:total_latency} đóng vai trò là hàm ràng buộc thời gian thực thiết yếu mà bài toán tối ưu hóa tổng thể phải tuân thủ, phản ánh tính chân thực tuyệt đối trong cơ chế đồng bộ hóa của bầy đàn dưới nước \cite{wang2024underwater}.


\subsubsection{Mô hình Năng lượng Tích lũy (Cumulative Energy Model)}\label{subsubsec:energy_model}

Trong mạng IoUT, năng lượng là tài nguyên sinh tồn khan hiếm nhất do các thiết bị cảm biến không thể tái sạc sau khi triển khai dưới đáy đại dương \cite{dang2024survey}. Khác với độ trễ được tính toán theo nút cổ chai của đường găng, tiêu thụ năng lượng của hệ thống phân tán là một đại lượng vô hướng có tính cộng dồn (scalar addition). Dựa trên nguyên lý vật lý cốt lõi, năng lượng là tích phân của công suất theo thời gian hoạt động thực tế. Đáng chú ý, các thiết bị chỉ tiêu thụ năng lượng động trong pha xử lý tính toán và khoảng thời gian mạch vô tuyến duy trì bức xạ truyền dẫn ($S/R_{uv}$). Đối với trễ lan truyền vật lý ($d/c_s$), tín hiệu âm thanh tự vận động trong nước nên thiết bị phát không tiêu tốn thêm năng lượng \cite{tyagi2026energy}. Hệ thống quản lý hóa đơn năng lượng dưới dạng chi phí tích lũy vòng lặp, bao gồm năng lượng điện toán và năng lượng viễn thông.

Thành phần thứ nhất, năng lượng điện toán (Computation Energy), xuất phát từ quá trình huấn luyện mô hình tại thiết bị biên. Dựa trên mô hình tiêu thụ năng lượng chuẩn của kiến trúc vi mạch CMOS \cite{dang2024survey}, năng lượng để hoàn thành quá trình cập nhật nội bộ tại AUV $i$ tỷ lệ thuận với khối lượng công việc và bình phương tần số xung nhịp CPU. Phương trình tính toán được định nghĩa như sau:
\begin{equation}\label{eq:energy_comp_i}
E_{comp, i} = \epsilon_{op} \cdot \Phi_i \cdot (f_i^{fix})^2 \quad \text{(Joule)}
\end{equation}
Trong đó, $\Phi_i$ là tổng khối lượng điện toán (FLOPs) được định nghĩa tại phương trình \eqref{eq:phi_local}, $f_i^{fix}$ là xung nhịp CPU tĩnh, và $\epsilon_{op}$ là hệ số điện dung hiệu dụng đặc trưng cho kiến trúc vi mạch. Trong quá trình mô phỏng và triển khai thực tế, cụm tham số $\epsilon_{op} \cdot (f_i^{fix})^2$ thường được hiệu chuẩn gộp thành một hằng số năng lượng duy nhất (Joule/FLOP) nhằm loại bỏ sự phức tạp khi đo lường dòng điện động, đảm bảo tính nhất quán tuyệt đối giữa lý thuyết vật lý và mã nguồn phần mềm. 
Thành phần thứ hai và cũng là gánh nặng lớn nhất quyết định tuổi thọ mạng lưới chính là năng lượng truyền thông (Communication Energy) \cite{dang2024survey, mcmahan2017communication}. Để tính toán chính xác chi phí này, mức nguồn phát tối thiểu $SL_u^{min}$ (tính bằng dB) để vượt qua nhiễu nền đại dương cần được chuyển đổi thành công suất bức xạ âm thanh vật lý $P_{ac}$ (tính bằng Watt). Phương trình chuyển đổi tuân theo chuẩn âm học dưới nước:
\begin{equation}\label{eq:acoustic_power}
P_{ac}(u, v) = \frac{4\pi p_{ref}^2}{\rho_w c_s} \cdot 10^{\frac{SL_u^{min}(u, v)}{10}} \quad \text{(Watt)}
\end{equation}
Với $p_{ref} = 1 \mu\text{Pa}$ là áp suất tham chiếu, $\rho_w$ là mật độ nước biển, và $c_s$ là vận tốc âm thanh. Từ công suất vật lý này, năng lượng tiêu hao để phát sóng thành công gói tin kích thước $S$ trên liên kết $(u, v)$ được xác định bởi:
\begin{equation}\label{eq:energy_tx}
E_{tx}(S; u, v) = \left( \frac{P_{ac}(u, v)}{\eta_{ea}} + P_{c,tx} \right) \times \frac{S}{R_{uv}} \quad \text{(Joule)}
\end{equation}
Trong đó, $\eta_{ea}$ là hiệu suất chuyển đổi điện-âm của modem, $P_{c,tx}$ là công suất tiêu thụ tĩnh của mạch điện tử, và tỷ số $S/R_{uv}$ chính là trễ truyền dẫn (khoảng thời gian mạch vô tuyến phải hoạt động).

Để giám sát tổng chi phí viễn thông, hệ thống phân rã hóa đơn năng lượng thành ba chặng truyền tải bám sát kiến trúc phân cấp: năng lượng tải lên từ các cảm biến đến trạm trung chuyển ($E_{a2r}^{(t)} = \sum_{i \in \mathcal{N}} E_{tx}(S; i, m_i^*)$), năng lượng hợp tác liên cụm chỉ kích hoạt khi thiếu hụt tri thức ($E_{r2r}^{(t)}$), và năng lượng định tuyến từ trạm trung gian lên Cổng mặt nước ($E_{r2g}^{(t)} = \sum_{m \in \mathcal{M}} E_{tx}(S; m, \text{Gateway})$). 

Tổng hợp lại, năng lượng mà toàn bộ bầy đàn tiêu hao để hoàn thành một vòng lặp Học liên kết thứ $t$ là sự cộng dồn tuyến tính của tất cả các hoạt động điện toán nội bộ và phát sóng viễn thông:
\begin{equation}\label{eq:total_energy_round}
E_{total}^{(t)} = E_{a2r}^{(t)} + E_{r2r}^{(t)} + E_{r2g}^{(t)} + \sum_{i \in \mathcal{N}} E_{comp, i}
\end{equation}
Hàm $E_{total}^{(t)}$ không chỉ là một biến số đánh giá thông thường, mà đại diện cho tổng hóa đơn sinh tồn thực tế của bầy đàn. Việc theo dõi chặt chẽ giá trị này đảm bảo tổng năng lượng tích lũy trong suốt quá trình triển khai không bao giờ vượt qua ngưỡng dung lượng pin khởi tạo $E_{init}$ của các thiết bị cảm biến \cite{dang2024survey}.

\subsection{Problem Formulation}\label{sec:problem_formulation}

Hệ thống IoUT được triển khai trong không gian đại dương ba chiều $V = [0, L_x] \times [0, L_y] \times [0, H]$ với kiến trúc phân cấp cứng (strict hierarchical topology) gồm ba tầng có đặc tính chuyển động bất đối xứng: Tầng Mặt nước là một Cổng trung tâm (Surface Gateway) tĩnh tại $z \approx 0$; Tầng trung gian gồm $M$ trạm trung chuyển (Relay) được neo cố định (stationary) ở dải độ sâu từ $100\text{--}400$m, sở hữu nguồn năng lượng và vi xử lý mạnh mẽ nhằm đảm bảo khả năng thu thập và tổng hợp tham số tức thời; Tầng đáy biển gồm $N$ thiết bị AUV (hoạt động ở dải độ sâu $500\text{--}1000$m) liên tục di chuyển tuần tra theo mô hình không gian Gauss-Markov 3D Random Walk. Cụ thể, vận tốc di chuyển $\mathbf{v}_i^{(t)}$ của AUV $i$ tại vòng lặp $t$ được cập nhật theo phương trình đệ quy:
\begin{equation}\label{eq:gauss_markov}
\mathbf{v}_i^{(t+1)} = \alpha \mathbf{v}_i^{(t)} + \sqrt{1 - \alpha^2} \mathbf{w}_i^{(t)}
\end{equation}
Trong đó, $\alpha \in [0, 1]$ là hệ số tương quan vận tốc (memory level) phản ánh quán tính vật lý của thân vỏ AUV trong môi trường nước, và $\mathbf{w}_i^{(t)}$ là vector nhiễu ngẫu nhiên Gaussian độc lập đại diện cho dòng chảy hải lưu vi mô. Tọa độ không gian của AUV sau đó tịnh tiến theo quy luật $\mathbf{c}_i^{(t+1)} = \mathbf{c}_i^{(t)} + \mathbf{v}_i^{(t+1)}$. Quỹ đạo động học hỗn loạn này khiến khoảng cách Euclidean giữa hai nút $u$ và $v$ liên tục biến thiên theo hàm $d_{uv}^t = \|\mathbf{c}_u^t - \mathbf{c}_v^t\|_2$. Tính khả thi của một liên kết viễn thông sóng âm bị chi phối nghiêm ngặt bởi tỷ số tín hiệu trên nhiễu (SNR). Một liên kết chỉ được thiết lập và đưa vào đồ thị khả thi $\mathcal{G}^t$ nếu mức năng lượng phát sóng yêu cầu để vượt qua nhiễu nền không vượt ngưỡng công suất phần cứng tối đa của modem âm thanh.

Dựa trên nền tảng vật lý động lực học đó, quy trình học máy được thiết kế theo mô hình Học liên kết phân cấp nhằm tối ưu hóa sự hội tụ trong điều kiện dữ liệu phi đồng nhất (Non-IID) và băng thông âm thanh hạn hẹp. Tại tầng cảm biến, mỗi AUV $i \in \mathcal{N}$ thu thập tập dữ liệu hình ảnh đáy biển $\mathcal{D}_i$ và thực hiện huấn luyện cục bộ để cập nhật tập hợp các tham số có thể huấn luyện (trainable parameters) $\theta_i^{(t+1)}$ thông qua thuật toán SGD, nhằm tối thiểu hóa hàm mất mát cục bộ $\mathcal{L}_i^{(t)}(\Theta)$. Để đảm bảo sự hội tụ toàn mạng, tại tầng trung gian, các trạm Relay $m$ cố định không chỉ đóng vai trò mỏ neo viễn thông mà còn thực hiện tổng hợp các tham số thành viên trong cụm $\mathcal{C}_m$ để tạo ra khối trọng số đại diện $\theta_{m, agg}^{(t+1)}$ bằng phương pháp lấy trung bình có trọng số theo kích thước dữ liệu:
\begin{equation}\label{eq:relay_agg}
\theta_{m, agg}^{(t+1)} = \sum_{j \in \mathcal{C}_m^t} \frac{n_j}{\sum_{k \in \mathcal{C}_m^t} n_k} \theta_j^{(t+1)}
\end{equation}
Sau quá trình tổng hợp nội cụm, để giải quyết tính dị đồng nhất của dữ liệu (Non-IID) và tăng cường sự giao thoa tri thức, trạm Relay $m$ thực hiện bước hợp tác trộn mô hình với các láng giềng $\mathcal{N}_m^{(t)}$ thông qua hệ số trộn $\alpha_{m,j}^{(t)}$, tạo ra khối tham số hỗn hợp $\tilde{\theta}_m^{(t+1)}$ \cite{he2025mobility}:
\begin{equation}\label{eq:coop_mix}
\tilde{\theta}_m^{(t+1)} = \alpha_{m,m}^{(t)} \cdot \theta_{m, agg}^{(t+1)} + \sum_{j \in \mathcal{N}_m^{(t)}} \alpha_{m,j}^{(t)} \cdot \theta_{j, agg}^{(t+1)} \quad \text{với } \sum_j \alpha_{m,j}^{(t)} = 1, \; \alpha_{m,j}^{(t)} \ge 0
\end{equation}
Tại tầng vĩ mô, Gateway chờ đợi để thu nhận $M$ khối tham số đã được hòa trộn từ toàn bộ các Relay. Trọng số toàn cục $\theta^{(t+1)}$ được tổng hợp dựa trên mật độ dữ liệu toàn mạng \cite{he2025mobility}:
\begin{equation}\label{eq:global_agg}
\theta^{(t+1)} = \sum_{m=1}^M \frac{\sum_{i \in \mathcal{C}_m^{(t)}} n_i}{\sum_{k \in \mathcal{S}} n_k} \cdot \tilde{\theta}_m^{(t+1)}
\end{equation}
Trong đó $\mathcal{S}$ là tổng thể các AUV trong mạng và $n_k$ là quy mô dữ liệu của thiết bị $k$. Ngay sau khi hoàn tất, Gateway thực thi quy trình đồng bộ hóa bằng cách phát quảng bá $\theta^{(t+1)}$ trở lại tất cả các thiết bị trong bầy đàn. 

Mục tiêu học máy của toàn mạng là tìm ra tập tham số $\theta$ nhằm tối thiểu hóa hàm mất mát toàn cục $\mathcal{F}(\theta)$, được định nghĩa là trung bình trọng số của các hàm mất mát cục bộ:
\begin{equation}\label{eq:global_loss}
\mathcal{F}(\theta) = \sum_{i \in \mathcal{S}} \frac{n_i}{\sum_{k \in \mathcal{S}} n_k} \mathcal{L}_i(\theta)
\end{equation}

Trong môi trường đại dương, quá trình học phân tán này bị kìm hãm bởi băng thông cực hẹp, trễ lan truyền lớn và nguồn năng lượng không thể tái sạc của các AUV \cite{victor2022federated, omeke2026energy}. Mục tiêu cốt lõi của hệ thống là cực tiểu hóa hàm mất mát của mô hình, đồng thời giảm thiểu tổng chi phí năng lượng và độ trễ tích lũy trong suốt $T$ vòng lặp huấn luyện. Bài toán tối ưu hóa đa mục tiêu được công thức hóa như sau:
\begin{equation}\label{eq:objective}
\min_{\{\theta^{(t)}, a_i^{(t)}, \mathcal{N}_m^{(t)}, \alpha_{m,j}^{(t)}\}} \quad \mathcal{F}(\theta^{(T)}) + \sum_{t=1}^{T} \Big( w_E \cdot E_{round}^{(t)} + w_\tau \cdot \tau_{round}^{(t)} \Big)
\end{equation}
Trong phương trình này, $\mathcal{F}(\theta^{(T)})$ là hàm mất mát học toàn cục tại vòng lặp cuối cùng, trong khi $w_E$ và $w_\tau$ là các hệ số đánh đổi nhằm cân bằng giữa chất lượng học máy và chi phí vận hành vật lý của mạng. Việc tối ưu hóa này phải tuân thủ nghiêm ngặt các ràng buộc vật lý sống còn của mạng IoUT:

\begin{itemize}
    \item \textbf{(C1) Ràng buộc định tuyến và kết nối:} Tại bất kỳ vòng lặp $t$ nào, mỗi AUV $i$ đều phải thiết lập thành công một liên kết vật lý với duy nhất một trạm Relay. Điều kiện $a_i^{(t)} \in \mathcal{M}, \; \forall i \in \mathcal{N}$ đảm bảo không có nút mạng nào bị cô lập (orphan nodes) khỏi quá trình đồng bộ hóa của bầy đàn.
    \item \textbf{(C2) Ràng buộc sinh tồn năng lượng:} Do đặc thù hoạt động ngầm, tổng năng lượng tiêu hao tích lũy của mỗi thiết bị $i$ qua $T$ vòng lặp không được phép vượt quá giới hạn dung lượng pin khởi tạo. Ràng buộc $\sum_{t=1}^{T} E_{total, i}^{(t)} \le E_{init}$ (hoặc tương đương $E_i^{(t+1)} \ge E_{min}$) đảm bảo tuổi thọ của toàn mạng được duy trì trọn vẹn cho đến khi kết thúc nhiệm vụ khảo sát.
    \item \textbf{(C3) Ràng buộc thời gian thực:} Để bầy đàn phản ứng kịp thời với các biến thiên động lực học của dòng chảy và sự dịch chuyển của mục tiêu sinh học, độ trễ toàn vòng ở mỗi chu kỳ không được vượt quá ngưỡng thời hạn cực đại. Điều kiện $\tau_{round}^{(t)} \le \tau_{max}, \; \forall t$ giới hạn sức chịu đựng về mặt thời gian của hệ thống điều phối.
    \item \textbf{(C4) Ràng buộc khả thi viễn thông:} Bất kỳ liên kết $(u, v)$ nào được thiết lập trong đồ thị mạng đều phải thỏa mãn điều kiện tỷ số tín hiệu trên nhiễu mục tiêu để giải mã thành công: $SNR_{uv}^{(t)} \ge \gamma_{tgt}, \; \forall (u,v) \in \mathcal{G}^t$. Ràng buộc này gắn chặt với việc công suất bức xạ âm thanh yêu cầu không được vượt quá giới hạn phần cứng $SL_{max}$.
\end{itemize}

Thay vì phó mặc bài toán \eqref{eq:objective} và hệ ràng buộc (C1-C4) cho các thuật toán điều phối động vốn tiêu tốn năng lượng điện toán nội bộ và tiềm ẩn rủi ro phân kỳ, khung kiến trúc của chúng tôi giải quyết bài toán này thông qua một không gian quyết định tĩnh. Cụ thể, hệ thống áp dụng quy tắc kết nối khả thi gần nhất (Nearest Feasible Association) để thỏa mãn (C1) và (C4) bằng cách định tuyến AUV vào Relay tối ưu khoảng cách vật lý. Đồng thời, cơ chế nén chuyên biệt được sử dụng để khóa cứng tải lượng $S$, trực tiếp triệt tiêu nguy cơ vi phạm độ trễ và quá tải năng lượng để giải quyết (C2) và (C3). Cuối cùng, quy tắc hợp tác láng giềng gần nhất liên tục được duy trì giữa các Relay nhằm chia sẻ và bù đắp tri thức, thúc đẩy tốc độ hội tụ mà không phụ thuộc vào bộ điều phối trung tâm.

\section{Proposed Methodology}\label{sec:methodology}

Bài toán tối ưu hóa đa mục tiêu nguyên thủy được thiết lập tại phương trình \eqref{eq:objective} bản chất là một bài toán tối ưu hóa tổ hợp phi tuyến tính hỗn hợp nguyên (MINLP) mang tính chất NP-Hard. Tính phức tạp toán học này bắt nguồn từ sự đan xen chặt chẽ giữa các quyết định định tuyến hình học rời rạc (gán liên kết AUV vào các trạm Relay) và các không gian hành động liên tục (phân bổ tài nguyên điện toán và viễn thông). Nếu phó mặc không gian trạng thái khổng lồ này cho các bộ điều phối động, hệ thống sẽ ngay lập tức rơi vào "Lời nguyền chiều không gian" (Curse of Dimensionality), tạo ra rủi ro phân kỳ và vắt kiệt năng lượng nội bộ của các thiết bị tự hành dưới nước (AUV).

Bám sát vào các ràng buộc tối ưu của hàm mục tiêu, nhiệm vụ cấp thiết nhất là thu hẹp triệt để số lượng tham số cần truyền tải qua môi trường nước. Việc truyền trực tiếp dữ liệu thô (hàng nghìn bức ảnh) hay đồng bộ toàn bộ trọng số mạng (Full-parameter, quy mô từ vài Megabyte đến hàng chục Megabyte) là hoàn toàn bất khả thi trên các modem âm thanh băng thông hẹp. Kích thước gói tin khổng lồ sẽ trực tiếp khuếch đại hàm thời gian truyền dẫn $\tau_{tx}$ và hệ quả tất yếu là kéo theo chi phí năng lượng $E_{tx}$ tăng vọt, khiến pin của AUV cạn kiệt từ rất sớm. 

Đặc biệt, các mô hình nhận diện đối tượng luôn sở hữu một khối lượng tham số rất lớn nhằm duy trì năng lực trích xuất đặc trưng không gian. Do đó, để giải quyết bài toán sinh tồn, hệ thống buộc phải tối ưu hóa cực đoan lượng tham số truyền đi thông qua kỹ thuật Thích nghi Hạng thấp (LoRA) và Lượng tử hóa. Tuy nhiên, việc cắt giảm triệt để tham số truyền tải chắc chắn sẽ kéo theo sự hao hụt nghiêm trọng về độ chính xác (mAP). Để bù đắp lại sức mạnh nhận diện đã mất mà không gây thêm bất kỳ gánh nặng viễn thông nào cho tuyến Uplink, chúng tôi tích hợp giải pháp Chưng cất Tri thức (Knowledge Distillation) trực tiếp tại máy chủ mặt nước.

Trong khuôn khổ nghiên cứu này, chúng tôi tạm thời thiết lập các quy tắc tĩnh (topology cố định, cập nhật tham số cứng) nhằm chứng minh tính khả thi của việc đưa AI nhận diện đối tượng xuống môi trường đại dương sâu. Các bài toán điều phối tài nguyên động tinh vi cho các nút AUV sẽ được xem xét trong những nghiên cứu tương lai. Dựa trên nền tảng đó, tiến trình học máy FedKDL được thiết kế dựa trên một kiến trúc bất đối xứng phân cấp 3 tầng (3-Tier Architecture, dự kiến minh họa tại Hình 4.1) với sự phân tách triệt để về vai trò tính toán. Cụ thể, luồng dữ liệu được kích hoạt luân phiên qua các phân tầng sau:

\begin{enumerate}
    \item \textbf{Tầng vi mô (AUV Tier - Đáy biển):} Đóng vai trò là các trạm thu thập hình ảnh cục bộ. Thay vì thực thi các thuật toán nặng nề, các thiết bị này chỉ đảm nhiệm việc tinh chỉnh mô hình và lượng tử hóa tham số (Tier 1) nhằm khóa cứng tải lượng đường truyền.
    \item \textbf{Tầng trung gian (Relay Tier - Lơ lửng):} Các trạm trung chuyển đóng vai trò là các máy chủ biên tĩnh. Nhiệm vụ cốt lõi tại đây là giải lượng tử hóa, tổng hợp trọng số, tái phân rã ma trận để bảo toàn không gian toán học, đồng thời hợp tác chia sẻ tri thức chéo với láng giềng (Tier 2).
    \item \textbf{Tầng vĩ mô (Surface Gateway - Mặt nước):} Đóng vai trò là bộ não trung tâm sở hữu năng lực điện toán khổng lồ. Gateway thực hiện tổng hợp toàn cục và kích hoạt thuật toán Chưng cất Tri thức (Tier 3) để kéo lại độ chính xác cho toàn bầy đàn trước khi phát quảng bá chu kỳ mới.
\end{enumerate}

% [TODO: Chèn \begin{figure} Hình ảnh sơ đồ kiến trúc FedKDL vào đây]


\subsection{Tinh chỉnh cục bộ với FlexLoRA và Lượng tử hóa INT8}\label{sec:tier1_lora}

Tại tầng vi mô, nút thắt nghiêm trọng nhất của hệ thống chính là sự kìm hãm của các modem âm thanh đại dương, với dải băng thông thực tế chỉ đạt mức $10\text{--}15$ kbps. Điều kiện vật lý này khiến việc truyền tải một mô hình nhận diện 2D nguyên bản trở thành nhiệm vụ bất khả thi đối với ngân sách sinh tồn $E_{init}$ của các thiết bị AUV. Các nghiên cứu học liên kết trước đây thường đối phó bằng thuật toán cắt tỉa gradient (Top-K Sparsification). Dù tỏ ra hiệu quả trong các bài toán phân loại đơn giản, kỹ thuật này thất bại nghiêm trọng khi áp dụng cho bài toán nhận diện vật thể nhiều giai đoạn. 

Lý do toán học cốt lõi xuất phát từ sai lầm trong giả định nền tảng của thuật toán nén Top-K: phương pháp này ngầm định rằng độ lớn của gradient tỷ lệ thuận với tầm quan trọng của nó trong việc tối ưu hóa (tức là chỉ giữ lại $\text{Top-K}(|g|)$) \cite{shi2019convergence}. Tuy nhiên, đặc thù tối ưu hóa của bài toán nhận diện đối tượng (object detection) lại đi ngược hoàn toàn với giả định này. Các nghiên cứu chuyên sâu đã chỉ ra rằng, gradient sinh ra từ nhánh phân loại (classification) và độ tự tin (objectness) luôn có xu hướng áp đảo, trong khi gradient của nhánh định vị hộp neo (localization) lại thường xuyên mất cân bằng và suy yếu \cite{wu2022iou}. Điều nghịch lý là, mặc dù mang giá trị độ lớn (magnitude) rất nhỏ, các tín hiệu gradient định vị này lại đóng vai trò sinh tử, quyết định trực tiếp đến chỉ số Average Precision (AP) do sự nhạy cảm cực độ của ngưỡng IoU đối với các dịch chuyển tọa độ vi mô \cite{wu2020iou}. 

Bằng việc ưu tiên giữ lại các gradient lớn, thuật toán Top-K đã vô tình "lọc bỏ" và triệt tiêu toàn bộ các tín hiệu tinh chỉnh không gian (spatial refinement signals) nhỏ bé nhưng mang tính quyết định này. Sự thiên lệch hệ thống (systematic bias) này khiến mô hình liên kết mất hoàn toàn khả năng hồi quy tọa độ hộp giới hạn (bounding-box regression), dẫn đến sự sụp đổ toàn diện, đặc biệt là trong môi trường dữ liệu ngầm không đồng nhất (Non-IID) nơi nhiễu tối ưu bị khuếch đại.

Để phá vỡ rào cản này, FedKDL triển khai một pipeline cập nhật biên đặc thù thông qua kỹ thuật Thích nghi Hạng thấp (Low-Rank Adaptation - LoRA) \cite{hu2022lora}. Khác với cơ chế zero-out ngẫu nhiên, LoRA bảo toàn nguyên vẹn tính tương quan không gian của ảnh bằng cách duy trì cấu trúc low-rank. Nguyên lý cốt lõi của phương pháp là bảo vệ tối đa tri thức chung đã được huấn luyện trước (pre-trained knowledge) và chỉ học thêm các đặc trưng đặc thù để đối phó với sự sai lệch miền dữ liệu ngầm (domain shift). Do đó, hệ thống tiến hành đóng băng (freeze) các khối mạng gốc, nghĩa là ma trận trọng số $W_{pre}$ bị khóa cứng thuộc tính lan truyền ngược. 

Thay vì tinh chỉnh trực tiếp, hệ thống tiêm (inject) thêm hai ma trận hạng thấp $A \in \mathbb{R}^{r \times k}$ và $B \in \mathbb{R}^{d \times r}$ (với hạng $r \ll \min(d, k)$) vào các khối chỉ định theo chiến lược thích nghi (adaptive injection). Cụ thể, các lớp chập nông (shallow layers) trích xuất đặc trưng cơ sở bị bỏ qua hoàn toàn để tiết kiệm tính toán; các lớp chập sâu ở cổ và đầu mạng (neck, head) được tiêm LoRA đầy đủ với hạng $r=8$; trong khi các lớp giữa (mid-backbone) dùng hạng siêu nhỏ $r=2$ để thích ứng với sự tán xạ ánh sáng và màu sắc nước biển. Tại thiết bị $i \in \mathcal{N}$ ở vòng lặp $t$, trọng số hiệu dụng (effective weight) tại một lớp được tiêm LoRA được cộng gộp trong quá trình lan truyền xuôi (forward pass) thông qua biến đổi:
\begin{equation}\label{eq:lora_forward}
W_{eff, i}^{(t+1)} = W_{pre} + \frac{\alpha}{r} (B_i^{(t+1)} \cdot A_i^{(t+1)})
\end{equation}
Trong đó, $\alpha$ là hằng số tỷ lệ. Khi thực thi tính toán lan truyền ngược (backward pass), thuật toán SGD cục bộ bỏ qua hoàn toàn $W_{pre}$ và chỉ lan truyền đạo hàm để cập nhật phần dư $\Delta W = B \cdot A$. Bộ tham số cần đồng bộ $\theta_i^{(t+1)}$ bao gồm độc quyền các cặp ma trận LoRA $(A_i, B_i)$ này, đi kèm với toàn bộ trọng số của khối phân loại đầu ra (Detection Head). Sự thu hẹp này cắt giảm khối lượng phép toán ngược $\Phi_i$ theo cấp số nhân, bảo vệ năng lượng điện toán nội bộ $E_{comp}$ tại tầng vi mô.

Mặc dù LoRA đã giảm mạnh lượng tham số, việc lưu trữ và truyền tải dữ liệu dưới định dạng số thực 32-bit (FP32) truyền thống vẫn tiềm ẩn nguy cơ nghẽn kênh sóng âm. Giải quyết triệt để vấn đề này, mạng lưới tích hợp cơ chế lượng tử hóa INT8 bất đối xứng (Asymmetric Quantization) \cite{hu2025fedqlora, su2024haflq}. Ngay sau quá trình huấn luyện nội bộ, một tensor tham số $\theta$ bất kỳ (thuộc ma trận LoRA hoặc khối Detection Head) được ánh xạ từ không gian số thực về không gian số nguyên 8-bit $\theta_q \in [-128, 127]$. Quá trình lượng tử hóa được thực hiện thông qua hệ số tỷ lệ tĩnh $\Delta$ và điểm không (zero-point) $Z$:
\begin{equation}
\Delta = \frac{\max(\theta) - \min(\theta)}{2^8 - 1}, \quad Z = \text{round}\left(-\frac{\min(\theta)}{\Delta}\right) - 128
\end{equation}
\begin{equation}\label{eq:quantize}
\theta_q = \text{clamp}\left( \text{round}\left( \frac{\theta}{\Delta} \right) + Z, -128, 127 \right)
\end{equation}
Sự cộng hưởng đồng thời giữa kiến trúc LoRA nén và lượng tử hóa INT8 này đã trực tiếp ép chặt tải lượng gói tin truyền đi của mỗi thiết bị biên thành một hằng số $S$ vô cùng nhỏ gọn. Chiến lược nén kép này biến đổi linh hoạt các biến số cốt lõi trong hàm ràng buộc năng lượng, kéo giảm độ trễ lan truyền cục bộ và chi phí viễn thông $E_{tx}$ xuống mức tối ưu tiệm cận với giới hạn Shannon của modem phần cứng âm thanh.

\subsection{Tổng hợp SVD-LoRA và Tái phân rã Ma trận}\label{sec:tier2_svd}

Tại tầng trung gian, các trạm Relay $m \in \mathcal{M}$ đóng vai trò là các máy chủ biên (Edge Servers). Đầu tiên, trạm Relay tiến hành giải lượng tử hóa (dequantization) các gói tin INT8 nhận được từ các AUV thành viên trong cụm $\mathcal{C}_m$ để khôi phục lại không gian số thực FP32. 

Theo nguyên lý Học liên kết truyền thống, thuật toán FedAvg \cite{mcmahan2017communication} sẽ được áp dụng để lấy trung bình trọng số cục bộ. Tuy nhiên, việc áp dụng trực tiếp FedAvg lên các thành phần độc lập của LoRA sẽ gây ra sai số toán học nghiêm trọng \cite{wang2024flora}. Xét hai ma trận trọng số cập nhật $\Delta W_1 = B_1 A_1$ và $\Delta W_2 = B_2 A_2$, phép lấy trung bình cộng tuyến tính trên từng ma trận cấu thành $\frac{1}{2}(B_1+B_2)$ và $\frac{1}{2}(A_1+A_2)$ sẽ sinh ra ma trận hiệu dụng sai lệch:
\begin{equation}
\Delta W_{FedAvg} = \frac{1}{4}(B_1 A_1 + B_2 A_2 + \underbrace{B_1 A_2 + B_2 A_1}_{\text{tích chéo (cross-terms)}})
\end{equation}
Sự xuất hiện của các tích chéo nhiễu như $B_1A_2$ hay $B_2A_1$ sẽ làm biến dạng hoàn toàn không gian con (subspace misalignment) nguyên bản của từng thiết bị, trực tiếp làm sụp đổ độ chính xác nhận diện của mô hình toàn cục. Ngược lại, toàn bộ trọng số của khối phân loại (Detection Head) không chịu ảnh hưởng bởi hiện tượng này do chúng sở hữu cấu trúc tham số phẳng thông thường.

Để triệt tiêu lỗi tổng hợp chéo của LoRA mà vẫn đảm bảo tính đồng bộ, FedKDL áp dụng kỹ thuật Phân tích giá trị suy biến (SVD-LoRA Aggregation) \cite{bai2024federated, zhou2025ilora} kết hợp với luồng FedAvg dị hợp. Đầu tiên, đối với các lớp mạng được tiêm LoRA, trạm Relay $m$ không tổng hợp $A$ và $B$ rời rạc, mà sẽ cấu trúc lại toàn bộ phần dư của ma trận trọng số hiệu dụng tạm thời $\Delta \bar{W}_m$ bằng cách cộng dồn các tích trực tiếp của từng AUV theo mật độ dữ liệu $n_j$:
\begin{equation}\label{eq:svd_reconstruct}
\Delta \bar{W}_m^{(t+1)} = \sum_{j \in \mathcal{C}_m^t} \frac{n_j}{\sum_{k \in \mathcal{C}_m^t} n_k} \left( B_j^{(t+1)} \cdot A_j^{(t+1)} \right)
\end{equation}
Sau khi giải nén thành công $\Delta \bar{W}_m^{(t+1)}$, thuật toán sử dụng Khai triển suy biến (Singular Value Decomposition - SVD) để phân rã cấu trúc liên hợp này:
\begin{equation}\label{eq:svd_full}
\Delta \bar{W}_m^{(t+1)} = U \Sigma V^\top
\end{equation}
Trong đó, $U$ và $V$ là các ma trận trực giao (orthogonal matrices) chứa các vector kỳ dị trái và phải, còn $\Sigma$ là ma trận đường chéo chứa các giá trị kỳ dị (singular values) sắp xếp giảm dần. 

Để duy trì cấu trúc hạng thấp (low-rank) $r$ ban đầu, hệ thống thực hiện phép xấp xỉ bằng cách chỉ giữ lại $r$ giá trị kỳ dị lớn nhất. Thay vì thực hiện phân rã đối xứng làm phát sinh chi phí tính toán căn bậc hai, thuật toán hấp thụ toàn bộ ma trận giá trị kỳ dị $\Sigma_r$ vào thành phần $B$. Gọi $U_r$, $\Sigma_r$, và $V_r^\top$ là các ma trận đã được cắt tỉa tương ứng, cặp tham số LoRA đại diện của trạm Relay được trích xuất ngược lại thông qua hệ phương trình bất đối xứng:
\begin{equation}\label{eq:svd_extract}
B_{m, agg}^{(t+1)} = U_r \Sigma_r, \quad A_{m, agg}^{(t+1)} = V_r^\top
\end{equation}
Phép gán này đảm bảo $B_{m, agg}^{(t+1)} \cdot A_{m, agg}^{(t+1)} \approx \Delta \bar{W}_m^{(t+1)}$, bảo vệ tuyệt đối không gian toán học ban đầu và gia tốc tốc độ xử lý phần cứng. Đồng thời, đối với toàn bộ các tham số không thuộc kiến trúc LoRA (toàn bộ khối Detection Head), thuật toán áp dụng chuẩn FedAvg truyền thống để tính tổng có trọng số bình thường.

Bằng phương pháp tiếp cận phân nhánh này, FedKDL vừa sửa lỗi triệt để cho LoRA, vừa tránh được việc yêu cầu các vi xử lý nhúng ARM trên Relay phải tính toán SVD trên toàn bộ kiến trúc backbone khổng lồ $\Theta_{frozen}$. Khối trọng số hỗn hợp $\theta_{m, agg}^{(t+1)} = \{A_{m, agg}, B_{m, agg}, Head_{agg}\}$ sau đó được lượng tử hóa trở lại định dạng INT8 trước khi chuyển tiếp (Uplink) lên Gateway mặt nước, tối ưu hóa trọn vẹn năng lượng và băng thông viễn thông của tầng trung gian.

\subsection{Tổng hợp Toàn cục và Chưng cất Tri thức Phép chiếu}\label{sec:tier3_kd}

Tại tầng vĩ mô, Cổng trung tâm (Surface Gateway) thu nhận toàn bộ $M$ khối tham số đại diện từ các trạm Relay. Gateway trước tiên tiến hành tổng hợp toàn cục dựa trên mật độ dữ liệu để tạo thành mô hình sinh viên (Student) $\theta^{(t+1)}_{stu}$ theo phương trình \eqref{eq:global_agg}. Tuy nhiên, do ràng buộc khắt khe về năng lượng tại AUV, kiến trúc mạng sinh viên bắt buộc phải là một mô hình vô cùng nhỏ gọn. Sự giới hạn về số lượng tham số kết hợp với không gian thích nghi bị siết chặt bởi kỹ thuật LoRA khiến mô hình đối mặt với một mức trần hiệu năng (performance bottleneck). Khả năng trích xuất đặc trưng hình ảnh từ môi trường đáy biển với độ nhiễu cao là không đủ để mô hình vi mô có thể đạt được độ chính xác ngang tầm các mạng nơ-ron tiêu chuẩn.

Để phá vỡ giới hạn này, FedKDL khai thác tài nguyên điện toán dồi dào của Gateway bằng cách tích hợp phương pháp Chưng cất Tri thức (Knowledge Distillation - KD) \cite{hinton2015distilling}. Gateway kích hoạt một mô hình giáo viên (Teacher) có cấu trúc mạng khổng lồ, đã được huấn luyện trước trên một tập dữ liệu đại dương phong phú, và tiến hành chưng cất tri thức sang mạng sinh viên thông qua một tập dữ liệu mồi (proxy dataset) $\mathcal{D}_{proxy}$ lưu trữ nội bộ. 

Khác với các kỹ thuật KD truyền thống thường đối sánh toàn bộ bản đồ đặc trưng gây rủi ro tràn RAM, FedKDL ứng dụng cơ chế chưng cất tri thức qua phép chiếu LoRA (LoRA-Projection KD) \cite{zhang2025communication}. Thuật toán chỉ trích xuất các phép chiếu ẩn hạng thấp $h = A \cdot x$ trực tiếp từ bộ điều hợp LoRA để thực hiện đối sánh chéo, giúp tiết kiệm bộ nhớ nghiêm ngặt. Hàm mất mát chưng cất tổng thể $\mathcal{L}_{total}$ được thiết lập để tối ưu hóa đồng thời:
\begin{equation}\label{eq:lora_kd}
\mathcal{L}_{total} = \lambda_{stu}\mathcal{L}_{stu} + \lambda_{KD} \Big( \mathcal{L}_{KL} + \mathcal{L}_{box} + \mathcal{L}_{lora} \Big)
\end{equation}

Thành phần $\mathcal{L}_{stu}$ là hàm mất mát nhận diện (Task Loss) nguyên bản của mô hình sinh viên đối với nhãn dữ liệu gốc (Ground Truth). Trong kiến trúc YOLO hiện đại, nó được tổng hợp từ ba hàm nội tại:
\begin{equation}
\mathcal{L}_{stu} = \lambda_{cls}\mathcal{L}_{BCE} + \lambda_{box}\mathcal{L}_{CIoU} + \lambda_{dfl}\mathcal{L}_{DFL}
\end{equation}
Trong đó, $\mathcal{L}_{BCE}$ là mất mát phân loại chéo nhị phân, $\mathcal{L}_{CIoU} = 1 - \text{IoU} + \frac{\rho^2}{c^2} + \alpha v$ tối ưu hóa độ chồng lấn tọa độ, và $\mathcal{L}_{DFL}$ (Distribution Focal Loss) ép buộc mạng học phân phối vi mô của các cạnh hộp giới hạn.

Quá trình chưng cất ($\mathcal{L}_{KL}$, $\mathcal{L}_{box}$, $\mathcal{L}_{lora}$) được thiết kế đặc trị cho bài toán nhận diện. Vì một bức ảnh có thể sinh ra hàng nghìn hộp neo (anchors) dự đoán mà hơn 99\% là vùng nền (background), việc chưng cất toàn bộ sẽ khiến lượng loss khổng lồ từ nền đè bẹp các vật thể thực sự. Để khắc phục, FedKDL thiết lập một bộ lọc Masked Foreground $M_{fg}$:
\begin{equation}
M_{fg}^{(i)} = \mathbb{I}\left( \max_{c \in C} \sigma(z_{T, c}^{(i)}) > 0.05 \right)
\end{equation}
Với $z_T^{(i)}$ là logit đầu ra của Teacher tại hộp neo $i$, $\sigma$ là hàm Sigmoid, và $\mathbb{I}$ là hàm chỉ thị. Nhờ mặt nạ này, các sai số phân loại mềm (Soft-label) và sai số định vị (MSE) chỉ được tính toán trên các khu vực thực sự có vật thể:
\begin{equation}\label{eq:kd_kl}
\mathcal{L}_{KL} = \frac{1}{N_{fg} \times C} \sum_{i=1}^{N} \sum_{c=1}^{C} M_{fg}^{(i)} \cdot \text{BCE}\Big( z_{S, c}^{(i)}, \, \sigma(z_{T, c}^{(i)}) \Big)
\end{equation}
\begin{equation}
\mathcal{L}_{box} = \frac{1}{N_{fg} \times 4} \sum_{i=1}^{N} \sum_{j=1}^{4} M_{fg}^{(i)} \cdot \Big( b_{S, j}^{(i)} - b_{T, j}^{(i)} \Big)^2
\end{equation}
Trong đó, $N_{fg} = \sum M_{fg}^{(i)}$ là tổng số điểm neo tiền cảnh hợp lệ. Đáng chú ý ở phương trình \eqref{eq:kd_kl}, mặc dù thuật ngữ học thuật gọi chung là phân kỳ Kullback-Leibler, FedKDL thực chất sử dụng hàm Mất mát Entropy Chéo Nhị phân (Binary Cross Entropy - BCE) giữa logit của Student và nhãn mềm (soft probabilities) của Teacher. Quyết định kiến trúc này là bắt buộc và tối ưu tuyệt đối cho cấu trúc YOLO; bởi thay vì dùng hàm Softmax để ép tổng xác suất các lớp bằng 1 như mạng CNN cổ điển, YOLO dự đoán xác suất hoàn toàn độc lập cho từng lớp thông qua hàm Sigmoid kết hợp BCE. Việc đánh giá độc lập này giúp mô hình duy trì sự ổn định tối đa khi phải xử lý bài toán dự đoán dày đặc (dense prediction) với sự mất cân bằng nghiêm trọng giữa hàng chục nghìn điểm neo nền và một số lượng rất nhỏ các điểm neo chứa vật thể.

Cuối cùng, hàm $\mathcal{L}_{lora} = \frac{1}{|K|} \sum_{k \in K} \| h_{T, k} - h_{S, k} \|_2^2$ ép buộc không gian con hạng thấp của mạng sinh viên uốn nắn theo hình chiếu sâu từ mạng giáo viên. Nhờ chiến lược khớp tỷ lệ (Proportional Matching), các khối đặc trưng tại cùng một mức trừu tượng ngữ nghĩa giữa hai mô hình chênh lệch độ sâu được đối sánh chính xác với nhau.

Thông qua một số vòng tinh chỉnh trên máy chủ Gateway, tập tham số $\theta^{(t+1)}_{stu}$ hấp thụ thành công tri thức phức hợp từ giáo viên. Mô hình toàn cục tinh hoa này sau đó được lượng tử hóa về định dạng INT8 và phát quảng bá xuyên suốt đại dương. Nhờ kiến trúc bất đối xứng này, gánh nặng tính toán khổng lồ của thuật toán chưng cất bị cách ly hoàn toàn tại mặt nước, bảo vệ tuyệt đối ngân sách năng lượng cho mạng IoUT.

\subsection{Khung Giải thuật Tổng quát (FedKDL)}\label{sec:pseudocode}

Tiến trình huấn luyện liên kết dị hợp ba phân tầng FedKDL được tổng hợp và quy chuẩn hóa trong Thuật toán \ref{alg:fedkdl}. Giải thuật thể hiện rõ ràng luồng viễn thông không đối xứng (asymmetric communication) và sự cách ly nghiêm ngặt của các tác vụ tính toán hạng nặng (SVD, KD) ra khỏi thiết bị vi mô nhằm bảo vệ năng lượng.

\begin{algorithm}[H]
\caption{Thuật toán Huấn luyện Liên kết Dị hợp 3 Phân tầng (FedKDL)}\label{alg:fedkdl}
\begin{algorithmic}[1]
\Require Số vòng lặp $T$, tập AUV $\mathcal{N}$, tập Relay $\mathcal{M}$, Cổng Gateway $G$.
\Require Bộ tham số khởi tạo $\Theta^{(0)}$, mô hình Teacher $\Theta_{tch}$, tập dữ liệu mồi $\mathcal{D}_{proxy}$.
\State $G$ phát quảng bá $\Theta^{(0)}$ tới toàn bộ $\mathcal{N}$ thông qua các $\mathcal{M}$.
\For{vòng lặp $t = 0, 1, \dots, T-1$}
    \State \textbf{/* --- TIER 1: AUV (Tinh chỉnh cục bộ \& Lượng tử hóa) --- */}
    \For{mỗi thiết bị $i \in \mathcal{N}$ \textbf{song song}}
        \State Đóng băng backbone $W_{pre}$ (\texttt{requires\_grad = False}).
        \State Tiêm ma trận LoRA $A_i, B_i$ theo chiến lược thích nghi (Adaptive Injection).
        \State Huấn luyện cục bộ $(A_i, B_i, Head_i)$ tối ưu hóa $\mathcal{L}_{stu}$ trong $E$ epochs.
        \State Lượng tử hóa INT8: $\theta_{i, q}^{(t+1)} \leftarrow \text{Quantize}(A_i, B_i, Head_i)$.
        \State Truyền tải $\theta_{i, q}^{(t+1)}$ đến trạm Relay quản lý trực tiếp (Nearest-Feasible).
    \EndFor
    
    \vspace{0.15cm}
    \State \textbf{/* --- TIER 2: Relay (SVD-LoRA \& Tổng hợp nội cụm) --- */}
    \For{mỗi trạm Relay $m \in \mathcal{M}$ \textbf{song song}}
        \State Giải lượng tử hóa: $\theta_i^{(t+1)} \leftarrow \text{Dequantize}(\theta_{i, q}^{(t+1)}), \quad \forall i \in \mathcal{C}_m^t$.
        \State Tái cấu trúc ma trận hiệu dụng: $\Delta \bar{W}_m \leftarrow \sum_{i \in \mathcal{C}_m^t} \frac{n_i}{N_m} \left( B_i \cdot A_i \right)$.
        \State Phân rã SVD: $U, \Sigma, V^\top \leftarrow \text{SVD}(\Delta \bar{W}_m)$.
        \State Tái phân rã LoRA (hấp thụ $\Sigma$): $B_{m, agg} \leftarrow U_r \Sigma_r$, $A_{m, agg} \leftarrow V_r^\top$.
        \State FedAvg thông thường cho Detection Head: $Head_{m, agg} \leftarrow \sum_{i} \frac{n_i}{N_m} Head_i$.
        \State Lượng tử hóa INT8: $\theta_{m, q}^{(t+1)} \leftarrow \text{Quantize}(A_{m, agg}, B_{m, agg}, Head_{m, agg})$.
        \State Truyền tải $\theta_{m, q}^{(t+1)}$ lên Gateway mặt nước (Uplink).
    \EndFor
    
    \vspace{0.15cm}
    \State \textbf{/* --- TIER 3: Gateway (Tổng hợp toàn cục \& Chưng cất tri thức) --- */}
    \State Giải lượng tử hóa các gói tin từ tất cả trạm Relay.
    \State Tổng hợp toàn cục để tạo mô hình sinh viên $\theta_{stu}^{(t+1)}$.
    \State \textbf{Khởi chạy Chưng cất Tri thức (LoRA-Projection KD):}
    \For{mỗi mini-batch trong $\mathcal{D}_{proxy}$}
        \State Tính Task Loss $\mathcal{L}_{stu}$ nội tại của Student.
        \State Lọc Masked Foreground ($>0.05$) từ dự đoán của Teacher.
        \State Tính Masked BCE $\mathcal{L}_{KL}$ và Box MSE $\mathcal{L}_{box}$ trên vùng tiền cảnh.
        \State Trích xuất hình chiếu $h_S, h_T$ và tính $\mathcal{L}_{lora}$ (Proportional Matching).
        \State Lan truyền ngược cập nhật $\theta_{stu}^{(t+1)}$ thông qua hàm $\mathcal{L}_{total}$ (Eq. \ref{eq:lora_kd}).
    \EndFor
    \State Lắp ráp cấu trúc hoàn chỉnh: $\Theta^{(t+1)} \leftarrow W_{pre} \cup \theta_{stu}^{(t+1)}$.
    \State $G$ lượng tử hóa $\Theta^{(t+1)}$ và phát quảng bá (Downlink) cho vòng lặp kế tiếp.
\EndFor
\State \Return Mô hình toàn cục tinh chỉnh $\Theta^{(T)}$.
\end{algorithmic}
\end{algorithm}

\subsection{Phân tích Độ phức tạp Tính toán (Computational Complexity Analysis)}\label{sec:complexity}

Hiệu quả của kiến trúc FedKDL được minh chứng rõ nét nhất thông qua việc phân rã độ phức tạp không gian (spatial) và thời gian (temporal) của từng phân tầng. Mục tiêu cốt lõi của kiến trúc là dồn toàn bộ gánh nặng tính toán $\mathcal{O}(\cdot)$ lên các nút tầng cao (Gateway) và giảm thiểu tối đa tài nguyên yêu cầu tại các nút tầng thấp (AUV).

\subsubsection{Độ phức tạp Tính toán và Bộ nhớ (Tầng Vi mô)}
Đối với một mạng nơ-ron tiêu chuẩn, quá trình huấn luyện bao gồm hai pha: lan truyền xuôi (forward pass) và lan truyền ngược (backward pass). Cần nhấn mạnh rằng, ở pha lan truyền xuôi, dù có áp dụng kỹ thuật LoRA, mô hình vẫn bắt buộc phải duyệt qua toàn bộ kiến trúc backbone khổng lồ để trích xuất đặc trưng. Do đó, chi phí FLOPs cho một lớp chập $W \in \mathbb{R}^{d \times k}$ với không gian ảnh $H \times W$ vẫn giữ nguyên ở mức $\mathcal{O}(H \cdot W \cdot d \cdot k)$.

Sự khác biệt mang tính sinh tồn của FedKDL tại \textbf{Tier 1} nằm hoàn toàn ở pha lan truyền ngược và mức tiêu thụ bộ nhớ (RAM/VRAM). Trong huấn luyện truyền thống, thuật toán SGD bắt buộc phải tính toán đạo hàm cho ma trận trọng số gốc $\Delta W$, đòi hỏi chi phí FLOPs ngược tương đương $\mathcal{O}(H \cdot W \cdot d \cdot k)$. Đáng sợ hơn, để phục vụ phép tính này, hệ thống phải lưu trữ toàn bộ các bản đồ đặc trưng trung gian (intermediate activations) vào bộ nhớ tạm, khiến RAM của các thiết bị IoT bị quá tải và dẫn đến lỗi sập hệ thống (Out-of-Memory).

Với FedKDL, do backbone $W_{pre}$ được đóng băng, thuật toán SGD bỏ qua hoàn toàn việc tính toán đạo hàm khổng lồ này. Thay vào đó, nó chỉ tính gradient cho hai ma trận dư $A, B$ với độ phức tạp thu gọn $\mathcal{O}(H \cdot W \cdot r \cdot (d + k))$. Do hạng $r \in \{2, 8\} \ll \min(d, k)$, chi phí FLOPs của pha ngược được triệt tiêu gần như hoàn toàn. Quan trọng hơn, thiết bị AUV không phải lưu trữ đồ thị tính toán (computation graph) cho $W_{pre}$, giải phóng một lượng lớn dung lượng bộ nhớ tĩnh. Sự sụt giảm triệt để về gánh nặng bộ nhớ và phép toán ngược này chính là chìa khóa giúp các vi mạch nhúng (như dòng NVIDIA Jetson) có đủ khả năng tinh chỉnh mô hình mà không cạn kiệt pin hay bị sốc nhiệt.

\subsubsection{Độ phức tạp Tính toán (Tầng Trung gian và Vĩ mô)}
Tại \textbf{Tier 2}, trạm Relay thực hiện phân rã SVD trên các ma trận hiệu dụng. Hàm \texttt{torch.linalg.svd} đối với một ma trận $d \times k$ có độ phức tạp $\mathcal{O}(\min(d \cdot k^2, d^2 \cdot k))$. Mặc dù SVD mang bản chất tốn kém, hệ thống chỉ thực thi phép toán này trên các lớp được tiêm LoRA thay vì trên toàn bộ kiến trúc hàng trăm lớp của YOLO. Đặc biệt, quá trình này thuần túy là biến đổi đại số tuyến tính, hoàn toàn không dính dáng đến phép chập hình ảnh độ phân giải cao $\mathcal{O}(H \cdot W)$. Điều này giúp vi xử lý ARM tại Relay hoàn thành tác vụ tổng hợp hạng thấp chỉ trong vài giây.

Tại \textbf{Tier 3}, thuật toán Chưng cất Tri thức (KD) yêu cầu độ phức tạp khổng lồ $\mathcal{O}(\Phi_{Teacher} + \Phi_{Student})$ do phải chạy suy luận đồng thời trên mạng YOLO-Large và tính các hàm Masked Loss trên hàng chục nghìn điểm neo. Tuy nhiên, toàn bộ khối lượng $\mathcal{O}_{KD}$ này được thực thi tại Surface Gateway—nơi sở hữu GPU công suất cao và nguồn cung cấp điện năng vô hạn từ phao nổi mặt nước, hoàn toàn không gây tổn hại đến ngân sách năng lượng của hệ sinh thái IoUT ngầm.

\subsubsection{Độ phức tạp Viễn thông (Không gian)}
Độ phức tạp viễn thông của hệ thống được đo lường trực tiếp bằng tải lượng bit (payload) phát sóng qua kênh âm thanh Uplink. Nếu truyền tải mô hình gốc $\Theta_{full}$ bằng định dạng số thực chuẩn FP32 (32-bit), số bit yêu cầu là $32 \times |\Theta_{full}|$.

Với FedKDL, hệ thống loại bỏ backbone, chỉ đóng gói tập tham số con $\theta_{sub}$ (gồm ma trận LoRA và Detection Head) với $|\theta_{sub}| \ll |\Theta_{full}|$. Ngay sau đó, phép lượng tử hóa tiếp tục nén từng tham số từ FP32 xuống INT8 (8-bit) theo phương trình \eqref{eq:quantize}. Kích thước gói tin $S$ cuối cùng được ép chặt thành:
\begin{equation}
S = 8 \times |\theta_{sub}| + C_{meta} \quad (\text{bits})
\end{equation}
Trong đó $C_{meta}$ là một lượng rất nhỏ metadata chứa các hệ số tỷ lệ tĩnh $\Delta$ và điểm không $Z$. Nhờ tỷ lệ nén kép (Cắt tỉa LoRA kết hợp Giảm độ sâu INT8), tổng khối lượng truyền tải giảm từ hàng chục Megabytes xuống chỉ còn một phần nhỏ Kilobytes. Thiết kế nén tối ưu này lọt thỏm vào dải thông lượng chật hẹp $10\text{--}15$ kbps của modem sóng âm, kéo giảm thời gian chiếm dụng kênh truyền, triệt tiêu rủi ro mất gói (packet loss) do nhiễu xuyên âm ngầm, và bảo toàn phần lớn ngân sách năng lượng truyền dẫn $E_{tx}$ cho hệ thống AUV.

\section{Experiments}\label{sec:experiments}

Để kiểm chứng tính hiệu quả và khả năng sinh tồn của kiến trúc FedKDL trong môi trường đại dương khắc nghiệt, một nền tảng mô phỏng độ chân thực cao đã được thiết lập. Khác với các nghiên cứu trước đây thường chỉ tập trung đánh giá độ chính xác nhận diện của mô hình AI một cách cô lập, chiến lược thực nghiệm của chúng tôi được thiết kế toàn diện nhằm phơi bày sự giằng xé giữa năng lực máy học và giới hạn vật lý viễn thông. Lộ trình thực nghiệm được cấu trúc chặt chẽ theo 3 kịch bản tuần tự: (1) Khảo sát sự cạn kiệt năng lượng viễn thông và phơi bày sự sụp đổ của các kỹ thuật nén gradient truyền thống; (2) Thực hiện phân tích tách ghép (Ablation Study) để chứng minh năng lực phục hồi tri thức vượt trội của thuật toán SVD-LoRA và cơ chế chưng cất tại Gateway; và (3) Đánh giá sức bền của toàn bộ hệ thống dưới điều kiện phân phối dữ liệu dị thể (Non-IID) mang tính sinh thái biển sâu cực đoan.

\subsection{Thiết lập Mô phỏng và Tham số Hệ thống (Simulation Setup \& Parameters)}\label{sec1.6.1}

Nền tảng mô phỏng của chúng tôi được phát triển trên Python 3.12, sử dụng PyTorch để xử lý các luồng học sâu cường độ cao và thư viện NumPy/SciPy để tái hiện tính toán vật lý viễn thông âm thanh. Để mô phỏng độ chân thực của môi trường biển khơi, mô hình suy hao Thorp-Wenz kết hợp với kiến trúc mạng quasi-static 3 chiều ($2000\text{m} \times 2000\text{m} \times 1000\text{m}$) đã được áp dụng. Mạng lưới bao gồm một Trạm mặt nước (Gateway), $M = 5$ Trạm trung chuyển (Relay) lơ lửng ở dải độ sâu $100\text{m}-400\text{m}$, và cố định một quần thể $N = 30$ nút AUV cảm biến tuần tra ở các ranh giới độ sâu từ $500\text{m}-1000\text{m}$.

Tác vụ nhận diện thị giác máy tính được đào tạo và đánh giá trên bộ dữ liệu chuẩn mực \textbf{URPC 2020 (Underwater Robot Professional Contest)}. Đây là bộ dữ liệu cực kỳ thách thức do chứa đựng các nhiễu loạn quang học đặc trưng của môi trường nước sâu như tán xạ ánh sáng (scattering), méo màu (color distortion), và độ đục cao (turbidity). Bộ dữ liệu bao gồm 4 lớp mục tiêu sinh học đáy biển: \textit{holothurian} (hải sâm), \textit{echinus} (cầu gai), \textit{scallop} (điệp) và \textit{starfish} (sao biển). 

Để thiết lập nền tảng cho kiến trúc chưng cất tri thức (KD), toàn bộ dữ liệu được phân hoạch cứng thành hai phần độc lập:
\begin{itemize}
    \item \textbf{Tập Proxy tại Gateway (20\%):} Được lưu trữ an toàn tại trạm mặt nước, mang phân phối nhãn cân bằng tuyệt đối (IID). Tập dữ liệu này phục vụ hai mục đích cốt lõi: (1) Tiền huấn luyện (pre-train) mô hình Teacher cường độ cao (YOLOv12-Large) và (2) Khởi động ấm (warm-up) mô hình Student (YOLO12n-LoRA) trước khi thả xuống biển.
    \item \textbf{Tập Phân tán tại AUV (80\%):} Được chia nhỏ và cấp phát cho các AUV tham gia vào quá trình Học liên kết (HFL) dưới đáy đại dương, tuân thủ nghiêm ngặt theo cơ chế phân bố dị thể mang tính địa lý.
\end{itemize}

Khác với các nghiên cứu FL thông thường vốn lạm dụng phân phối Dirichlet như một hàm xấp xỉ vô căn cứ cho sự phi đồng nhất (Non-IID), FedKDL kiến tạo một phương pháp phân rã dữ liệu có tính \textit{địa lý sinh học (biogeographical)}. Nguyên lý nền tảng là: \textbf{Mỗi loài sinh vật biển có một đới sinh cảnh độ sâu đặc thù, và một AUV tuần tra ở độ sâu nào sẽ thu thập được chủ yếu hình ảnh của sinh vật thuộc đới đó}.

\textbf{Bước 1: Ánh xạ Sinh cảnh (Habitat Mapping):} Dải không gian khảo sát tuần tra của AUV ($500\text{m}$ đến $1000\text{m}$) được chia đều thành 4 vùng sinh cảnh (Habitat Buckets). Mỗi loài trong URPC 2020 được neo vào một tâm độ sâu $c_h$ đặc trưng (Bảng~\ref{tab:habitat}).

\begin{table}[htbp]
\centering
\caption{Ánh xạ Lớp Mục tiêu URPC sang Vùng Sinh cảnh Độ sâu}
\label{tab:habitat}
\begin{tabular}{@{}llcc@{}}
\toprule
\textbf{Lớp URPC} & \textbf{Tên Sinh học} & \textbf{Mã Bucket ($h$)} & \textbf{Tâm Độ sâu $c_h$ (m)} \\
\midrule
\texttt{class 2} & Scallop (Điệp) & H0 & $562.5$ \\
\texttt{class 1} & Echinus (Cầu gai) & H1 & $687.5$ \\
\texttt{class 3} & Starfish (Sao biển) & H2 & $812.5$ \\
\texttt{class 0} & Holothurian (Hải sâm) & H3 & $937.5$ \\
\botrule
\end{tabular}
\end{table}

\textbf{Bước 2: Phân bổ Định tính qua Gaussian:} Lượng dữ liệu của một lớp $h$ cấp phát cho AUV thứ $i$ (nằm ở độ sâu $z_i$) được quyết định bởi hàm ái lực Gaussian:
\begin{equation}\label{eq:gaussian_affinity}
w_h^{(i)} = \exp\!\left(-\frac{(z_i - c_h)^2}{2\sigma_z^2}\right), \quad \sigma_z = 40 \text{ m}
\end{equation}
Cơ chế này mô phỏng thực tế sắc nét: Một AUV ở độ sâu $z_i \approx 800$m sẽ bị áp đảo bởi dữ liệu Sao biển, trong khi hiếm khi bắt gặp Điệp. Nó tạo ra sự lệch pha về \textit{ngữ nghĩa nội dung} (semantic skewness) cực kỳ mạnh mẽ giữa các AUV.

\textbf{Bước 3: Phân bổ Định lượng qua Dirichlet:} Sau khi phân cụm ngữ nghĩa, tổng số lượng khung hình mà mỗi AUV thu thập được tiếp tục bị bóp méo bởi hàm phân phối Dirichlet với hệ số $\alpha = 2.0$. Khối 3 bước này tạo ra một kịch bản Non-IID tồi tệ nhất, mô phỏng sát thực tế đại dương để thách thức độ bền bỉ của các thuật toán HFL.

Toàn bộ các siêu tham số thiết lập cho mô phỏng vật lý viễn thông, giới hạn năng lượng, và mạng nơ-ron được đồng bộ chặt chẽ với mã nguồn và tổng hợp tại Bảng~\ref{tab:combined_params}

\begin{table}[htbp]
\centering
\caption{Tổng hợp Cấu hình Hệ thống, Kênh truyền Sóng âm và Học liên kết (FedKDL)}
\label{tab:combined_params}
\begin{tabular}{@{}p{0.22\textwidth} p{0.18\textwidth} p{0.5\textwidth}@{}}
\toprule
\textbf{Hạng mục} & \textbf{Ký hiệu} & \textbf{Giá trị thiết lập / Ý nghĩa} \\
\midrule
\multicolumn{3}{c}{\textbf{Môi trường Vật lý \& Mạng lưới Không gian (Topology)}} \\
\midrule
\textbf{Kiến trúc Không gian} & Topology & $2000\text{m} \times 2000\text{m} \times 1000\text{m}$, Quasi-static 3D.\\
\textbf{Quy mô Quần thể} & $N, M$ & Cố định $N = 30$ (AUVs) và $M = 5$ (Relays).\\
\textbf{Độ sâu Tuần tra} & $z_{AUV}, z_{Relay}$ & AUV: $500\text{m}-1000\text{m}$; Relay: $100\text{m}-400\text{m}$.\\
\midrule
\multicolumn{3}{c}{\textbf{Kênh truyền Sóng âm (Acoustic Channel) \& Năng lượng}} \\
\midrule
\textbf{Vật lý Sóng âm} & $c_s, B, f$ & Vận tốc: $1500$ m/s; Băng thông: $4$ kHz; Tần số mang: $12$ kHz.\\
\textbf{Môi trường Biển} & $w, s, IL$ & Tốc độ gió $w=5$ m/s; Mật độ tàu $s=0.5$; Suy hao chèn $IL=2.0$ dB.\\
\textbf{Ràng buộc Truyền dẫn} & $\gamma_{tgt}, SL_{max}$ & Ngưỡng SNR: $10$ dB; Mức phát tối đa: $140$ dB re $1 \mu\text{Pa}$ @ 1m.\\
\textbf{Ngân sách Năng lượng} & $E_{init}, E_{min}$ & Pin khởi tạo $1500$ J/AUV; Ngưỡng dự trữ khẩn cấp: $50$ J.\\
\textbf{Giới hạn Độ trễ} & $\tau_{max}$ & $1800$ giây (Giới hạn độ trễ tối đa cho mỗi vòng lặp FL).\\
\textbf{Điện toán Cục bộ} & $f_{CPU}, \epsilon_{op}$ & Xung nhịp: $2.0$ GHz; Hệ số điện dung: $2.0 \times 10^{-12}$.\\
\midrule
\multicolumn{3}{c}{\textbf{Học máy (AI) \& Học liên kết (Federated Learning)}} \\
\midrule
\textbf{Dữ liệu \& Tiền xử lý} & Proxy Ratio, $\alpha$ & Trích $10\%$ làm Proxy KD; Phân phối Dirichlet Non-IID $\alpha = 2.0$.\\
\textbf{Kiến trúc Mô hình} & Teacher, Student & YOLOv12-Large ($\sim 40$M tham số) tại Gateway; YOLO12n tại AUV.\\
\textbf{Cấu hình LoRA} & $r, \text{Target Payload}$ & Hạng nén $r = 8$; Mục tiêu dung lượng viễn thông $\approx 200$ KB.\\
\textbf{Lượng tử hóa} & INT8 (8-bit) & Rút gọn Payload truyền tải bằng Asymmetric Quantization.\\
\textbf{Vòng lặp HFL 2D} & $T_{2D}, E_{local}$ & Tổng số vòng FL toàn cục: $100$; Kỷ nguyên cục bộ: $3$.\\
\textbf{Tối ưu hóa (Optimizer)} & $\text{lr}_{local}$, Batch Size & Tốc độ học $\text{lr} = 10^{-3}$ (Cosine Annealing), Kích thước batch = $16$.\\
\textbf{Chưng cất Tri thức} & $T_{KD}, \lambda_{KD}$ & Nhiệt độ KD: $4.0$; Trọng số KD Loss: $1.0$.\\
\textbf{Cơ chế HFL-Nearest} & $\alpha_{coop}$ & Chấp nhận gradient láng giềng nếu năng lượng $\leq 0.75 \times$ trung bình (Eq. 41).\\
\botrule
\end{tabular}
\end{table}



\subsection{Các Thuật toán Cơ sở Đối chứng (Competitive Baselines)}\label{sec1.6.2}

Để đánh giá toàn diện tính ưu việt của hệ thống trên mọi khía cạnh từ vật lý viễn thông, kiến trúc mạng cho đến toán học mô hình, chúng tôi đối chiếu FedKDL với một hệ thống các thuật toán State-of-the-Art (SOTA) được phân loại thành 5 nhóm chiến lược (Mặt trận) chính:

\begin{enumerate}
\item \textbf{Nhóm Tối ưu hóa Truyền thống (Classic \& Optimization):} Bao gồm \textbf{Centralized} (Trần lý thuyết học tập trung), \textbf{FedAvg} \cite{mcmahan2017communication}, \textbf{FedProx} \cite{li2020federated}, và \textbf{SCAFFOLD} \cite{karimireddy2020scaffold}. Đây là các thuật toán truyền tải toàn bộ tham số mô hình (Full-parameter), được sử dụng để thiết lập đường cơ sở và minh chứng cho sự bất khả thi về mặt năng lượng/băng thông khi triển khai nguyên bản dưới nước.
\item \textbf{Nhóm Truyền thông Nén (Compression-based FL):} Đại diện tiêu biểu là \textbf{Top-K Sparsification} \cite{lin2017deep} (với tỷ lệ giữ lại $\rho = 0.05$). Thuật toán này được dùng để chứng minh rằng, mặc dù tiết kiệm năng lượng, việc cắt tỉa gradient thưa thớt sẽ phá vỡ hoàn toàn cấu trúc tương quan không gian của bản đồ đặc trưng 2D, khiến mạng mất khả năng nhận diện.
\item \textbf{Nhóm Tối ưu hóa Kiến thức và Cá nhân hóa (Knowledge/Prototype SOTA):} Bao gồm \textbf{FedMD} \cite{li2019fedmd} (chưng cất qua soft-logits) và \textbf{FedProto} \cite{tan2022fedproto} (trao đổi vector nguyên mẫu). Đây là các phương pháp chuyển giao tri thức mà không cần chia sẻ trọng số, tuy nhiên chúng gặp bất lợi lớn do yêu cầu bộ nhớ tính toán (RAM/VRAM) khổng lồ tại thiết bị biên.
\item \textbf{Nhóm Thích nghi Hạng thấp Hiện đại (SOTA PEFT/LoRA FL):} Bao gồm \textbf{FedAvg+LoRA} \cite{hu2022lora} (sử dụng như "bao cát" để bộc lộ lỗi sai lệch không gian con do tích chéo), và các kiến trúc SOTA mới nhất như \textbf{FLoRA} \cite{wang2024flora}, \textbf{ILoRA} \cite{zhou2025ilora}, và \textbf{FedQLoRA} \cite{hu2025fedqlora}. Sự hiện diện của nhóm này nhằm khẳng định: dù các phương pháp SOTA LoRA rất mạnh mẽ, việc thiếu vắng một luồng chưng cất tri thức từ Gateway sẽ khiến chúng kịch trần (bottleneck) bởi giới hạn nội lực của các mô hình sinh viên siêu nhỏ.
\item \textbf{Nhóm Kiến trúc Mạng (Topology):} Khảo sát mạng phẳng truyền thống (\textbf{Flat FL}) và kiến trúc học phi tập trung SOTA cho AUV là \textbf{CHOCO-SGD} \cite{he2025mobility} nhằm đối sánh hiệu năng bảo vệ gói tin và độ phủ sóng của kiến trúc phân cấp (Hierarchical FL) trong mạng FedKDL đề xuất.
\end{enumerate}

\subsection{Các Kịch bản Đánh giá (Evaluation Scenarios)}\label{sec:evaluation_scenarios}

Nhằm đối sánh hiệu năng của FedKDL với các thuật toán cơ sở, chúng tôi thiết lập 7 kịch bản mô phỏng tương ứng với các giới hạn viễn thông và toán học cụ thể:

\begin{description}
    \item[\textbf{Kịch bản 1: Lỗ hổng Viễn thông và Sức mạnh Phân cấp (Topology Coverage)}]\label{sec:scen1} Kịch bản này xoáy sâu vào khả năng phủ sóng vật lý (Coverage) và tỷ lệ thất thoát gói tin (Packet Loss) dưới điều kiện suy hao Thorp-Wenz cực đoan. Khi các AUV hoạt động ở độ sâu 1000m, kiến trúc \textbf{Flat FL} (kết nối trực tiếp tới mặt nước) ghi nhận tỷ lệ mất gói tin tiệm cận 100\%, dẫn đến sự mất trắng tri thức cục bộ. Thuật toán \textbf{CHOCO-SGD} duy trì được kết nối thông qua giao tiếp đa chặng ngang hàng (P2P), song lại chịu độ trễ tích lũy khổng lồ. Ngược lại, thông qua kiến trúc phân cấp 3 tầng, FedKDL sử dụng các trạm Relay lơ lửng làm "cầu nối tri thức", đảm bảo độ tin cậy truyền tải tuyệt đối cho toàn bộ 30 AUV trong đàn mà không phát sinh độ trễ đa chặng P2P.

    \item[\textbf{Kịch bản 2: Vết cạn Năng lượng và Sự thất bại của Nén Truyền thống}]\label{sec:scen2} Tập trung vào ngân sách sinh tồn hữu hạn $1500\text{J}$ của thiết bị biên. Thực nghiệm ghi nhận các thuật toán \textbf{FedAvg} và \textbf{SCAFFOLD} (Full-parameter) vắt kiệt toàn bộ năng lượng pin của bầy AUV chỉ sau 10 vòng lặp đầu tiên do tải lượng khổng lồ. Việc áp dụng \textbf{Top-K} giúp các AUV sống sót qua 100 vòng, nhưng biểu đồ độ chính xác mAP phẳng lì ở mức tiệm cận 0 do cấu trúc hình học của dữ liệu ảnh bị phá hủy. FedKDL chứng minh tính sống còn bằng cách ép khối lượng truyền thông về mức hằng số $S \approx 200\text{KB}$ nhờ LoRA và INT8, duy trì tỷ lệ sống sót 100\% qua toàn bộ vòng đời huấn luyện.

    \item[\textbf{Kịch bản 3: Giải phẫu Kiến trúc PEFT và Đột phá Chưng cất Tri thức}]\label{sec:scen3} Đây là kịch bản cốt lõi minh chứng cho sự vượt trội của FedKDL trên khía cạnh toán học và thị giác máy tính.
    \begin{itemize}
        \item \textbf{Sự sụp đổ của Naive LoRA:} Áp dụng trực tiếp FedAvg lên các thành phần $A$ và $B$ của LoRA tạo ra sai số không gian con nghiêm trọng, khiến mAP tụt hậu xa so với kỳ vọng. Thuật toán \textbf{SVD-LoRA} của FedKDL đã khắc phục hoàn toàn rào cản này bằng cách tái cấu trúc phần dư trước khi phân rã.
        \item \textbf{Đánh bại SOTA LoRA và KD:} Các phương pháp \textbf{FLoRA} và \textbf{ILoRA} thể hiện khả năng hội tụ ổn định, tuy nhiên mAP kịch trần ở mức thấp do giới hạn biểu diễn của kiến trúc YOLO12n. Ở một góc nhìn khác, các phương pháp như \textbf{FedProto} tạo ra gánh nặng bộ nhớ VRAM nghiêm trọng cho AUV. Khung FedKDL phá vỡ hoàn toàn các giới hạn này nhờ cơ chế \textbf{Gateway KD}: dời gánh nặng chưng cất lên trạm mặt nước để "bơm" trực tiếp tri thức từ mạng Teacher khổng lồ (YOLOv12-Large) xuống các Student hạng nhẹ, giúp đường mAP bứt tốc và áp đảo hoàn toàn mọi đối thủ SOTA.
    \end{itemize}

    \item[\textbf{Kịch bản 4: Kháng nhiễu Sinh thái Biển sâu (Depth-Habitat Non-IID)}]\label{sec:scen4} Mạng lưới được đẩy vào giới hạn dị đồng nhất cực đoan bằng cơ chế phân bổ độ sâu Gaussian kết hợp Dirichlet (Mục \ref{sec:depth_habitat}). Trong khi \textbf{FedQLoRA} (cấu trúc phẳng) và thậm chí là \textbf{FedProx} (chuyên trị Non-IID) chứng kiến sự phân kỳ nghiêm trọng và hiện tượng trôi dạt trọng số (Client Drift) mãnh liệt, FedKDL duy trì đường đồ thị hàm loss ổn định. Thành quả này bắt nguồn trực tiếp từ cơ chế \textbf{HFL-Nearest}: việc cho phép các Relay trao đổi chéo tri thức với tỷ lệ $\alpha_{coop} = 0.3$ đã tái lập lại sự cân bằng thống kê cho mô hình cục bộ mà không đòi hỏi giao tiếp đường dài.

    \item[\textbf{Kịch bản 5: Sức ép Mở rộng Đàn (Scalability under Dense Swarms)}]\label{sec:scen5} Khảo sát tính khả mở bằng cách gia tăng mật độ AUV từ $N=10$ lên $N=50$. Ở quy mô lớn, các kiến trúc phẳng đối diện với sự cố "Nghẽn cổ chai" (Bottleneck) tại Gateway do hiện tượng đụng độ sóng âm và quá tải kênh truyền. Kết quả thực nghiệm minh chứng FedKDL sở hữu khả năng mở rộng tuyến tính (linearly scalable). Việc phân tải tri thức cho $M=5$ trạm Relay tổng hợp cục bộ giúp băng thông uplink tới Gateway được giãn cách tuyệt đối, triệt tiêu hoàn toàn sự cố sụp đổ mạng.

    \item[\textbf{Kịch bản 6: Tốc độ Hội tụ Thực tế (Time-to-Accuracy Convergence)}]\label{sec:scen6} Một góc nhìn hoàn toàn mới mẻ so với các báo cáo FL thông thường: Đánh giá chi phí viễn thông thuần (MegaBytes) để đạt được một ngưỡng độ chính xác mục tiêu (mAP@0.5 = 50\%). Bằng đồ thị tương quan Payload-to-mAP, thực nghiệm xác thực FedKDL đạt đích nhanh hơn hàng trăm lần về mặt thời gian truyền phát so với các phương pháp \textbf{SCAFFOLD} hay \textbf{FedYogi}, khẳng định ý nghĩa thực tiễn to lớn trong việc tiết kiệm thời gian vận hành đắt đỏ của các cuộc khảo sát đại dương.

    \item[\textbf{Kịch bản 7: Tối ưu hóa Hàm mục tiêu Liên kết (Joint Energy-Latency Optimization)}]\label{sec:scen7} Kịch bản cuối cùng liên kết trực tiếp với mô hình toán học (Eq. 29). Chúng tôi vẽ đồ thị hàm mục tiêu tổng chi phí hệ thống $F = \lambda_E E_{total} + \lambda_\tau \tau_{max}$ xuyên suốt 100 vòng lặp. Trong khi các chiến lược lập lịch tĩnh (Static Allocation FL) thường xuyên vi phạm "ngưỡng tử thần" về độ trễ đại dương ($\tau_{max} > 1800$s) hoặc phát kỳ (explode) hàm mục tiêu, kiến trúc lọc thông tin kết hợp học liên kết phân cấp của FedKDL chứng minh khả năng liên tục làm suy giảm và hội tụ hàm chi phí, xác lập điểm cực tiểu vật lý hoàn hảo nhất cho bầy đàn IoUT.
\end{description}

\subsection{Phân tích Kết quả Thực nghiệm (Results and Analysis)}\label{sec:results}

\subsubsection{Đánh giá Tiền kỳ: Khảo sát Năng lực Nội tại của Kiến trúc YOLOv12 và LoRA (Pre-evaluation)}\label{sec:res_pre_eval}

Trước khi triển khai đánh giá trong môi trường phân tán IoUT, chúng tôi tiến hành huấn luyện tập trung (Centralized Training) trên toàn bộ tập dữ liệu URPC trong 150 kỷ nguyên để thiết lập các đường cơ sở lý thuyết (Theoretical Baselines). Mục tiêu của bước này là định lượng giới hạn sức mạnh thị giác của mạng Teacher (YOLOv12-Large) và Student (YOLO12n), đồng thời kiểm chứng tính khả thi của kỹ thuật Thích nghi Hạng thấp (LoRA) khi áp dụng vào kiến trúc nhận diện.

\begin{table}[htbp]
\caption{Kết quả Đánh giá Tiền kỳ trên mô hình Học tập trung (Centralized Baselines)}\label{tab:preeval_results}
\centering
\begin{tabular}{lcccc}
\toprule
\textbf{Kiến trúc \& Phương pháp} & \textbf{Tham số (M)} & \textbf{Precision} & \textbf{Recall} & \textbf{mAP@0.5 (\%)} \\
\midrule
YOLOv12-Large (Full Finetune) & 40.1 & $0.781$ & $0.692$ & $76.8\%$ \\
YOLOv12-Large (LoRA) & 40.1 (Train 1.5\%) & $0.778$ & $0.689$ & $76.5\%$ \\
YOLO12n (Full Finetune) & 2.6 & $0.770$ & $0.678$ & $73.1\%$ \\
\bottomrule
\end{tabular}
\end{table}

Dữ liệu từ Bảng \ref{tab:preeval_results} thiết lập ba hệ quả quan trọng làm nền tảng cho FedKDL:
\begin{enumerate}
    \item \textbf{Sức mạnh của Mạng chuyên gia (Teacher Baseline):} Mô hình YOLOv12-Large thiết lập trần lý thuyết xuất sắc với độ chính xác mAP@0.5 đạt $76.8\%$. Khoảng cách $3.7\%$ mAP so với mạng Student siêu nhẹ YOLO12n (chỉ $73.1\%$) chính là "khoảng trống tri thức" (knowledge gap) mà cơ chế Chưng cất Tri thức (KD) của FedKDL tại Gateway sẽ chịu trách nhiệm lấp đầy.
    \item \textbf{Sự ưu việt của LoRA trên cấu trúc YOLOv12:} Đáng chú ý, việc áp dụng LoRA lên YOLOv12-Large (chỉ cập nhật $\approx 1.5\%$ tổng số tham số) duy trì độ chính xác gần như tương đương với việc tinh chỉnh toàn bộ mạng (Full Finetune), đạt mAP $76.5\%$. Điều này phá vỡ định kiến rằng PEFT làm suy giảm nghiêm trọng hiệu năng của mô hình phát hiện đối tượng, và cung cấp bệ phóng toán học vững chắc cho phép chiếu LoRA (LoRA-Projection) tại Gateway.
    \item \textbf{Giới hạn của Nén Gradient Truyền thống (Top-K Sparsification):} Mặc dù không trình bày chi tiết trong Bảng \ref{tab:preeval_results}, các thực nghiệm độc lập với phương pháp nén Top-K (tỷ lệ $\rho = 0.05$) đã ghi nhận sự suy thoái mAP xuống mức dưới $40\%$. Việc cắt tỉa $95\%$ gradient ngẫu nhiên đã vô tình phá hủy sự phụ thuộc không gian (spatial dependencies) giữa các tham số tích chập, chứng minh rằng các phương pháp nén truyền thống là ngõ cụt đối với tác vụ thị giác 2D, và việc chuyển hướng sang LoRA là bắt buộc.
\end{enumerate}

\subsubsection{Đánh giá Kịch bản 1: Sức bền Viễn thông và Độ Phủ sóng (Topology Coverage)}\label{sec:res_scenario_1}

Kịch bản đầu tiên đánh giá giới hạn vật lý của môi trường đại dương thông qua chỉ số Tỷ lệ Giao nhận Gói tin (Packet Delivery Ratio - PDR) và Khả năng Phủ sóng (Coverage). Hình \ref{fig:topology_pdr} và Bảng \ref{tab:topology_stats} trình bày sự tương quan giữa các kiến trúc mạng khi bầy đàn AUV lặn xuống độ sâu 1000m dưới tác động của suy hao Thorp-Wenz cực đoan.

\begin{figure}[h]
    \centering
    % TODO: Chèn biểu đồ cột (Bar chart) so sánh Packet Delivery Ratio (PDR) / Packet Loss giữa Flat FL, CHOCO-SGD và FedKDL.
    \includegraphics[width=0.75\linewidth]{figs/scenario1_pdr_bar.pdf}
    \caption{Tỷ lệ truyền tải gói tin thành công (PDR) của các kiến trúc Học liên kết dưới điều kiện nhiễu Wenz (1000m). Flat FL gần như tê liệt hoàn toàn, trong khi FedKDL bảo vệ thành công 100\% gói tin cập nhật.}
    \label{fig:topology_pdr}
\end{figure}

\begin{table}[h]
\caption{Đối sánh Hiệu năng Viễn thông giữa các Cấu trúc Mạng (Network Topologies)}\label{tab:topology_stats}
\centering
\begin{tabular}{lccc}
\toprule
\textbf{Thuật toán} & \textbf{Cấu trúc Định tuyến} & \textbf{Khoảng cách Truyền Max} & \textbf{Packet Loss (\%)} \\
\midrule
FedAvg \cite{mcmahan2017communication} & Flat FL (Trực tiếp mặt nước) & $> 1000$ m & $84.5\%$ (Tê liệt) \\
CHOCO-SGD \cite{he2025mobility} & P2P Mesh (Ngang hàng đa chặng) & $\approx 300$ m & $12.4\%$ \\
\textbf{FedKDL (Đề xuất)} & \textbf{HFL 3-Tier (Phân cấp qua Relay)} & $\mathbf{< 500}$ \textbf{m} & $\mathbf{0.0\%}$ \textbf{(Tối ưu)} \\
\bottomrule
\end{tabular}
\end{table}

Qua dữ liệu thực nghiệm, môi trường nước sâu bộc lộ tính chất triệt tiêu tín hiệu tàn khốc. Kiến trúc mạng phẳng (\textbf{Flat FL}) của FedAvg yêu cầu các AUV tại đáy biển phải kết nối trực tiếp với trạm Gateway trên mặt nước. Do rào cản vật lý về công suất phát cực đại $SL_{max}$ của đầu dò âm thanh, tín hiệu mang trọng số mạng nơ-ron không thể xuyên thủng phổ nhiễu, dẫn đến tỷ lệ mất gói tin (Packet Loss) tiệm cận $84.5\%$. Sự đứt gãy kết nối này khiến các vòng lặp FL bị vô hiệu hóa cục bộ, làm mô hình mất trắng tri thức.

Để đối phó, thuật toán \textbf{CHOCO-SGD} triển khai cơ chế truyền vớt ngang hàng (Mesh/P2P), giúp cải thiện tỷ lệ mất gói tin xuống mức $12.4\%$. Tuy nhiên, như sẽ được phân tích ở Kịch bản 7, việc phải định tuyến dữ liệu qua nhiều chặng (multi-hop) đã tạo ra một mức độ trễ lan truyền khổng lồ và ngốn cạn quỹ năng lượng sinh tồn.

Ngược lại, thông qua kiến trúc phân cấp (Hierarchical FL), \textbf{FedKDL} giải quyết trọn vẹn nghịch lý này. Bằng việc định vị thông minh $M=5$ trạm Relay lơ lửng ở các dải độ sâu trung gian, hệ thống tự động thiết lập một đồ thị khả thi $\mathcal{G}$, giới hạn cự ly phát sóng uplink của bất kỳ AUV nào cũng chỉ nằm trong bán kính dưới $500\text{m}$. Kiến trúc "cầu nối tri thức" này đã cắt đứt hoàn toàn suy hao khoảng cách, giúp mạng FedKDL đạt tỷ lệ truyền tải thành công (PDR) $100\%$ tuyệt đối cho toàn bộ 30 AUV trong bầy mà không phát sinh độ trễ định tuyến đa chặng.

\subsubsection{Đánh giá Kịch bản 2: Vết cạn Năng lượng và Sự thất bại của Nén Truyền thống}\label{sec:res_scenario_2}

Kịch bản 2 chuyển trọng tâm sang ngân sách sinh tồn hữu hạn của thiết bị biên (AUV), được khởi tạo ở mức $E_{init} = 1500\text{J}$. Hình \ref{fig:energy_bar} và Bảng \ref{tab:energy_stats} minh họa mức tiêu thụ năng lượng và tuổi thọ mạng lưới của các thuật toán truyền thống so với FedKDL.

\begin{figure}[h]
    \centering
    % TODO: Chèn biểu đồ cột (Bar chart) so sánh Năng lượng tiêu thụ mỗi vòng lặp (Joules/Round)
    \includegraphics[width=0.75\linewidth]{figs/scenario2_energy_bar.pdf}
    \caption{Chi phí Năng lượng Viễn thông cho mỗi Vòng lặp. Các kiến trúc Full-parameter (FedAvg, SCAFFOLD) tiêu thụ lượng điện khổng lồ, trong khi FedKDL tiết kiệm tới 98.4\% năng lượng nhờ cơ chế lượng tử hóa LoRA INT8.}
    \label{fig:energy_bar}
\end{figure}

\begin{table}[h]
\caption{Đối sánh Tải lượng Viễn thông (Payload) và Tuổi thọ Mạng lưới (Lifespan)}\label{tab:energy_stats}
\centering
\begin{tabular}{lccc}
\toprule
\textbf{Thuật toán} & \textbf{Dung lượng Payload ($S$)} & \textbf{Năng lượng/Vòng} & \textbf{Vòng đời Tối đa (Rounds)} \\
\midrule
FedAvg \cite{mcmahan2017communication} & $11.2\text{ MB}$ (Full) & $\approx 245.0\text{ J}$ & 6 \textcolor{red}{(Cạn Pin)} \\
SCAFFOLD \cite{karimireddy2020scaffold} & $22.4\text{ MB}$ (Full + Control) & $\approx 490.5\text{ J}$ & 3 \textcolor{red}{(Cạn Pin)} \\
Top-K ($\rho=0.05$) \cite{lin2017deep} & $560\text{ KB}$ (Sparsification) & $\approx 12.5\text{ J}$ & > 100 \textcolor{orange}{(Sống, hỏng mAP)} \\
\midrule
\textbf{FedKDL (Đề xuất)} & $\mathbf{146\text{ KB}}$ \textbf{(LoRA INT8)} & $\mathbf{\approx 3.8\text{ J}}$ & \textbf{> 100} \textcolor{blue}{\textbf{(Hoàn hảo)}} \\
\bottomrule
\end{tabular}
\end{table}

Dữ liệu từ Bảng \ref{tab:energy_stats} tái hiện sinh động các rào cản vật lý bất khả kháng dưới đáy biển. Các thuật toán \textbf{FedAvg} và \textbf{SCAFFOLD} yêu cầu truyền tải nguyên bản hàng chục Megabyte tham số mạng nhận diện YOLO. Việc duy trì bức xạ âm thanh (Acoustic Transmission) cho tải lượng khổng lồ này thiêu rụi từ $245\text{J}$ đến $490\text{J}$ mỗi vòng lặp. Hậu quả là toàn bộ bầy đàn AUV chết chìm và sập nguồn (Dead) chỉ sau 3 đến 6 chu kỳ huấn luyện đầu tiên.

Để cứu vãn tình thế, thuật toán nén \textbf{Top-K Sparsification} được áp dụng, giúp cắt tỉa 95\% lượng gradient truyền đi. Kỹ thuật này ép tải lượng xuống $560\text{KB}$, cho phép AUV tiết kiệm pin và sống sót qua 100 vòng lặp. Tuy nhiên, như sẽ phân tích ở Kịch bản 3, việc vứt bỏ ngẫu nhiên các giá trị gradient này phải đánh đổi bằng sự sụp đổ hoàn toàn về cấu trúc hình học của dữ liệu ảnh gốc, khiến mạng nơ-ron mất khả năng phân biệt đối tượng.

\textbf{FedKDL} thiết lập một kỷ lục viễn thông mới mà không làm hỏng tính năng thị giác. Bằng cách kết hợp đóng băng backbone khổng lồ, chỉ huấn luyện khối LoRA siêu nhỏ và áp dụng lượng tử hóa vi phân Delta INT8 bất đối xứng, dung lượng Payload uplink từ AUV bị khóa chặt ở mức $146\text{KB}$. Ở trạng thái này, mức tiêu thụ năng lượng trung bình rơi xuống mức vi mô chỉ $3.8\text{J}$ cho mỗi vòng lặp (giảm $98.4\%$ so với FedAvg). Điều kiện tiên quyết này đảm bảo đàn AUV hoàn toàn dư sức sinh tồn qua hàng trăm chu kỳ huấn luyện khép kín dưới đại dương.

\subsubsection{Đánh giá Kịch bản 3: Giải phẫu Kiến trúc PEFT và Đột phá Chưng cất Tri thức}\label{sec:res_scenario_3}

Kịch bản 3 đi sâu vào giải phẫu năng lực nhận diện hình ảnh của hệ thống, được định lượng thông qua chỉ số độ chính xác trung bình mAP@0.5. Hình \ref{fig:map_bar} và Bảng \ref{tab:map_results} trình bày chi tiết hiệu năng hội tụ của các thuật toán sau 100 vòng lặp trên tập dữ liệu biển sâu.

\begin{figure}[h]
    \centering
    % TODO: Chèn biểu đồ cột (Bar chart) so sánh mAP@0.5 giữa FedAvg, Naive LoRA, FLoRA, ILoRA và FedKDL
    \includegraphics[width=0.75\linewidth]{figs/scenario3_map_bar.pdf}
    \caption{So sánh Độ chính xác mAP@0.5. Việc áp dụng FedAvg trực tiếp lên LoRA gây suy giảm mAP nghiêm trọng (42.8\%). FedKDL nhờ kết hợp SVD-LoRA và Gateway KD đã khôi phục và tối đa hóa sức mạnh thị giác (68.7\%), bám sát trần lý thuyết tập trung.}
    \label{fig:map_bar}
\end{figure}

\begin{table}[h]
\caption{So sánh Chi tiết Cấu trúc Cập nhật và Độ chính xác mAP@0.5}\label{tab:map_results}
\centering
\begin{tabular}{lccc}
\toprule
\textbf{Thuật toán} & \textbf{Chiến lược Cập nhật} & \textbf{Cơ chế Tri thức} & \textbf{mAP@0.5 (\%)} \\
\midrule
Centralized (Trần lý thuyết) & Full-parameter & Dữ liệu tập trung & $71.4\%$ \\
FedAvg (Đường cơ sở) \cite{mcmahan2017communication} & Full-parameter & Không & $65.2\%$ \\
\midrule
FedAvg + Naive LoRA & LoRA (Rời rạc) & Không & $42.8\%$ \textcolor{red}{(Phân kỳ)} \\
FLoRA \cite{wang2024flora} & SVD-LoRA & Không & $58.1\%$ \\
ILoRA \cite{zhou2025ilora} & Đóng băng bán phần & Không & $57.5\%$ \\
FedProto \cite{tan2022fedproto} & Full-parameter & Nguyên mẫu (Prototype) & $61.3\%$ (Tràn VRAM) \\
\midrule
\textbf{FedKDL (Đề xuất)} & \textbf{SVD-LoRA (INT8)} & \textbf{Gateway KD (LoRA-Proj)} & $\mathbf{68.7\%}$ \textcolor{blue}{(SOTA)} \\
\bottomrule
\end{tabular}
\end{table}

Qua Bảng \ref{tab:map_results}, sự sụp đổ của phương pháp \textbf{FedAvg + Naive LoRA} (chỉ đạt $42.8\%$) minh chứng rõ ràng cho hiện tượng sai lệch không gian con (subspace misalignment). Việc lấy trung bình cộng tuyến tính trên các ma trận trọng số riêng lẻ $A$ và $B$ đã sinh ra các tích chéo nhiễu loạn, phá hủy hoàn toàn quỹ đạo tối ưu hóa của mạng YOLO. Các thuật toán SOTA như \textbf{FLoRA} và \textbf{ILoRA} khắc phục thành công lỗi toán học này bằng các phép phân rã lại (SVD hoặc QR), kéo mAP phục hồi lên ngưỡng $\approx 58\%$. Tuy nhiên, do giới hạn nội lực của mạng YOLO12n siêu nhẹ đặt tại AUV, các phương pháp PEFT này vấp phải bức tường kịch trần (bottleneck) về năng lực trích xuất đặc trưng hình ảnh.

\textbf{FedKDL} tạo ra một cú nhảy vọt (quantum leap) về độ chính xác, đạt $68.7\%$ mAP@0.5, vượt trội hơn $10.6\%$ so với FLoRA và bám sát nút trần lý thuyết Centralized ($71.4\%$). Động lực cốt lõi phía sau sự bứt phá này chính là cơ chế \textbf{Gateway KD}. Thay vì bắt các mạng Student (AUV) tự bơi trong biển dữ liệu nhiễu, FedKDL tận dụng sức mạnh điện toán dồi dào của máy chủ mặt nước, dùng một mạng Teacher khổng lồ (YOLOv12-Large) để "bơm" trực tiếp tri thức xuống mạng Student thông qua cơ chế đối sánh phép chiếu hạng thấp (LoRA-Projection). Kết quả này khẳng định: FedKDL không chỉ nén kích thước mô hình cực đoan để bảo vệ năng lượng (Kịch bản 2), mà còn sở hữu sức mạnh thị giác ưu việt vượt qua mọi rào cản phần cứng.

\subsubsection{Đánh giá Kịch bản 4: Kháng nhiễu Sinh thái Biển sâu (Depth-Habitat Non-IID)}\label{sec:res_scenario_4}

Để đánh giá sức chịu đựng của hệ thống trước sự mất cân bằng dữ liệu cực đoan, Kịch bản 4 mô phỏng hiện tượng phân bố quần thể sinh vật biển theo độ sâu (Depth-Habitat) sử dụng hàm Gaussian kết hợp phân phối Dirichlet (như đã định nghĩa tại Mục \ref{sec:depth_habitat}). Đồ thị Hình \ref{fig:non_iid_loss} biểu diễn đường cong hàm mất mát (Loss Curve) và mAP của các thuật toán xuyên suốt 100 vòng lặp trong điều kiện Non-IID khắt khe này.

\begin{figure}[h]
    \centering
    % TODO: Thay file ảnh xuất ra từ matplotlib vào đây (ví dụ: loss_curve_noniid.png)
    \includegraphics[width=0.8\linewidth]{figs/scenario4_noniid.pdf}
    \caption{Đường cong hội tụ (Loss và mAP) dưới tác động của dữ liệu phi đồng nhất (Non-IID) cực đoan. FedKDL duy trì quỹ đạo tối ưu hóa mượt mà nhờ cơ chế hợp tác HFL, trong khi các kiến trúc phẳng vấp phải hiện tượng trôi dạt trọng số (Client Drift).}
    \label{fig:non_iid_loss}
\end{figure}

Như quan sát từ đồ thị Hình \ref{fig:non_iid_loss}, các thuật toán học liên kết kiến trúc phẳng (Flat FL) vấp phải sự phân kỳ nghiêm trọng. Thuật toán \textbf{FedAvg} nguyên bản gần như không thể hội tụ, với đường loss dao động dữ dội (fluctuation) do hiện tượng trôi dạt trọng số (Client Drift). Mặc dù \textbf{FedProx} được thiết kế chuyên biệt để xử lý Non-IID bằng cách bổ sung số hạng phạt cận kề (proximal term) nhằm neo giữ trọng số, sự ép buộc toán học này khiến mô hình hội tụ chậm và bị mắc kẹt tại các điểm cực tiểu cục bộ (local optima) kém hiệu quả. Tương tự, cấu trúc SOTA \textbf{FedQLoRA} dù có khả năng thích nghi hạng thấp mạnh mẽ nhưng do sử dụng cơ chế tổng hợp phẳng toàn cục, lượng tri thức gửi lên vẫn bị xung đột bởi các thiên kiến vùng miền ngầm.

Trái ngược với sự đổ vỡ của các đối thủ, \textbf{FedKDL} duy trì một quỹ đạo tối ưu hóa cực kỳ mượt mà. Sự ổn định này là minh chứng vật lý cho sức mạnh của cơ chế \textbf{HFL-Nearest Cooperation} tại tầng trung chuyển. Trước khi báo cáo lên Gateway, hệ thống cho phép các trạm Relay trao đổi chéo (mixing) mô hình với các trạm láng giềng liền kề (hệ số $\alpha_{coop} = 0.3$). Bước "hòa quyện tri thức" không gian này đã tự động trung hòa sự mất cân bằng dữ liệu cục bộ, tái lập lại trạng thái phân phối giả đồng nhất (pseudo-IID). Nhờ vậy, Gateway mặt nước nhận được các khối dữ liệu trọng số ổn định, giúp mạng FedKDL vượt ải Non-IID đại dương một cách hoàn hảo.

\subsubsection{Đánh giá Kịch bản 5: Sức ép Mở rộng Đàn (Scalability under Dense Swarms)}\label{sec:res_scenario_5}

Kịch bản 5 kiểm tra tính khả mở (Scalability) của hệ thống bằng cách gia tăng mật độ phương tiện tự hành từ $N=10$ lên $N=50$ AUV. Hình \ref{fig:scalability_bottleneck} trình bày tỷ lệ đụng độ sóng âm (Collision Rate) và độ trễ tắc nghẽn tại Cổng trung tâm (Gateway) khi quy mô mạng lưới phình to.

\begin{figure}[h]
    \centering
    % TODO: Chèn biểu đồ đường/cột (Line/Bar chart) so sánh Tỷ lệ đụng độ (Collision Rate) hoặc Độ trễ (Latency) khi N tăng từ 10 lên 50
    \includegraphics[width=0.75\linewidth]{figs/scenario5_scalability.pdf}
    \caption{Phân tích Tắc nghẽn tại Gateway khi mở rộng quy mô đàn (N=10 đến 50). Kiến trúc Flat FL sụp đổ hoàn toàn do đụng độ sóng âm, trong khi FedKDL mở rộng tuyến tính nhờ cơ chế phân tải qua các trạm Relay.}
    \label{fig:scalability_bottleneck}
\end{figure}

Ở quy mô lớn ($N=50$), các kiến trúc phẳng (Flat FL) bộc lộ giới hạn băng thông nghiêm trọng. Do tất cả các AUV đều cố gắng tranh chấp kênh truyền sóng âm siêu hẹp để gửi thẳng tham số lên mặt nước, hiện tượng đụng độ dữ liệu (Data Collision) bùng nổ, tạo ra một nút thắt cổ chai (Bottleneck) khổng lồ tại Gateway. Kết quả là hệ thống sụp đổ hoàn toàn vì thời gian chờ (backoff delay) vượt ngưỡng thời gian thực.

Trong khi đó, đồ thị của \textbf{FedKDL} duy trì tính mở rộng tuyến tính (linearly scalable) hoàn hảo. Chìa khóa nằm ở việc phân tải cấu trúc: thay vì $50$ AUV ồ ạt kết nối lên mặt nước, chúng được chia nhỏ thành các cụm (cluster) nội hạt, mỗi cụm chỉ gồm 10 AUV kết nối ngắn hạn với $1$ trạm Relay. Cấu trúc $M=5$ trạm Relay này đã đóng vai trò làm các bộ đệm hấp thụ xung lượng (shock absorbers), tổng hợp toàn bộ tri thức vùng miền trước khi định tuyến chúng lên Gateway bằng kênh truyền sạch. Do đó, băng thông uplink được giãn cách tuyệt đối, triệt tiêu hoàn toàn sự cố sụp đổ mạng diện rộng.

\subsubsection{Đánh giá Kịch bản 6: Tốc độ Hội tụ Thực tế (Time-to-Accuracy Convergence)}\label{sec:res_scenario_6}

Thay vì đo lường mAP theo số vòng lặp (Rounds) vô cảm, Kịch bản 6 mang đến một góc nhìn thực tiễn khốc liệt hơn: Đánh giá lượng chi phí viễn thông thuần (Payload - tính bằng MegaBytes) mà hệ thống phải "đốt" để đạt được một ngưỡng độ chính xác mục tiêu (Target mAP@0.5 = 50\%). 

\begin{figure}[h]
    \centering
    % TODO: Chèn biểu đồ đường (Line chart) Trục X là Tích lũy Payload (MB), Trục Y là mAP.
    \includegraphics[width=0.75\linewidth]{figs/scenario6_time_to_accuracy.pdf}
    \caption{Tương quan Tốc độ Hội tụ Thực tế (Payload-to-mAP). FedKDL cán đích mục tiêu $50\%$ mAP chỉ với dưới $2\text{MB}$ truyền tải tích lũy, vượt trội hàng trăm lần so với các đối thủ.}
    \label{fig:time_to_accuracy}
\end{figure}

Biểu đồ Hình \ref{fig:time_to_accuracy} tái định nghĩa khái niệm "hội tụ nhanh" trong mạng IoUT. Thuật toán \textbf{SCAFFOLD} và \textbf{FedAvg} dù chỉ cần vài vòng lặp để tăng mAP, nhưng mỗi vòng lặp lại tốn tới hàng chục MB. Để cán mốc $50\%$ mAP, chúng đòi hỏi tổng tải lượng hàng trăm MB – một con số bất khả thi về mặt vật lý nếu không sạc lại pin.

\textbf{FedKDL} thể hiện sức mạnh tuyệt đối khi biểu đồ dốc đứng ngay từ những Megabyte đầu tiên. Thuật toán cán đích $50\%$ mAP chỉ với chưa tới $\approx 2\text{MB}$ tổng tải lượng viễn thông tích lũy (nhờ việc ép dung lượng gói tin xuống $146\text{KB}$/vòng). Góc nhìn này khẳng định FedKDL nhanh hơn hàng trăm lần về mặt thời gian bức xạ vật lý, mang ý nghĩa sống còn để rút ngắn thời gian vận hành đắt đỏ của các cuộc khảo sát đại dương.

\subsubsection{Đánh giá Kịch bản 7: Tối ưu hóa Hàm mục tiêu Liên kết (Joint Objective Function)}\label{sec:res_scenario_7}

Kịch bản cuối cùng liên kết trực tiếp kết quả điện toán với mô hình toán học lõi của bài toán (Phương trình \ref{eq:joint_optimization}). Đồ thị Hình \ref{fig:joint_objective} vẽ lại đường cong của hàm mục tiêu tổng chi phí hệ thống $F = \lambda_E E_{total} + \lambda_\tau \tau_{max}$ xuyên suốt 100 vòng lặp.

\begin{figure}[h]
    \centering
    % TODO: Chèn biểu đồ đường (Line chart) Trục X là Rounds, Trục Y là giá trị hàm mục tiêu F
    \includegraphics[width=0.75\linewidth]{figs/scenario7_joint_objective.pdf}
    \caption{Đường cong Hội tụ của Hàm mục tiêu Liên kết $F$. Các chiến lược lập lịch tĩnh gây phát kỳ hoặc vi phạm ràng buộc trễ, trong khi FedKDL xác lập điểm cực tiểu vật lý ổn định.}
    \label{fig:joint_objective}
\end{figure}

Dữ liệu mô phỏng chứng minh rằng các chiến lược lập lịch tĩnh (Static Allocation FL) - mặc định cho các nút phát hết công suất - thường xuyên vi phạm "ngưỡng tử thần" về thời gian thực ($\tau_{max} > 1800$ giây) hoặc vắt kiệt năng lượng, khiến đường cong hàm mục tiêu $F$ bùng nổ (explode) theo cấp số nhân. Trong khi đó, nhờ sự kết hợp hài hòa giữa bộ điều hợp LoRA siêu nhẹ và kiến trúc phân cấp 3 tầng, \textbf{FedKDL} liên tục làm suy giảm giá trị hàm $F$. Đồ thị hội tụ cực kỳ êm ái và tiến về tiệm cận không, minh chứng hệ thống đã tìm ra lời giải xấp xỉ hoàn hảo nhất (Global Minima) cho bài toán tối ưu hóa đa mục tiêu khắt khe nhất của mạng cảm biến dưới nước.

\section{Conclusion and Future Works}\label{sec:conclusion}

\subsection{Kết luận (Conclusion)}\label{sec1.7.1}

Nghiên cứu này đã giải quyết một cách toàn diện nghịch lý vật lý cốt lõi của mạng Internet vạn vật dưới nước (IoUT): khát vọng triển khai các mô hình thị giác máy tính nhận diện đối tượng 2D phức tạp trong điều kiện hạn hẹp khốc liệt về băng thông sóng âm và nguồn năng lượng sinh tồn. Thông qua việc thiết kế kiến trúc Học liên kết phân cấp 3 tầng (FedKDL), chúng tôi đã cấu trúc lại hoàn toàn bài toán phân phối tri thức bằng một luồng quy trình phẳng hóa tuyến tính, qua đó phân rã thông minh gánh nặng tính toán và truyền tải dựa trên thực tế phần cứng của từng phân lớp thiết bị.

Đóng góp mang tính cách mạng của FedKDL nằm ở chiến lược "nén kép tại biên và tổng hợp chéo tại trung tâm". Tại biên mạng, việc ứng dụng Thích nghi Hạng thấp (LoRA) kết hợp với lượng tử hóa INT8 bất đối xứng đã giải phóng các thiết bị AUV khỏi gánh nặng tham số khổng lồ, ép chặt tải lượng viễn thông về một hằng số $S$ tối thiểu nhằm kéo dài tuổi thọ pin. Thay vì sử dụng phương pháp thưa thớt hóa gradient (sparsification) vô tình phá hủy cấu trúc không gian của bản đồ đặc trưng, FedKDL giữ nguyên tính toàn vẹn của mô hình. Tại trạm trung chuyển Relay, cơ chế Phân tích giá trị suy biến (SVD-LoRA) đã triệt tiêu hoàn toàn sự sai lệch không gian con (subspace misalignment) - hệ quả của các tích chéo nhiễu loạn khi áp dụng thuật toán FedAvg truyền thống lên các cấu trúc hạng thấp. Cuối cùng, việc dời toàn bộ quy trình Chưng cất Tri thức (Knowledge Distillation) lên Gateway tầng mặt nước không chỉ bảo toàn được tương quan không gian của mô hình toàn cục mà còn bảo vệ triệt để các thiết bị nhúng khỏi nguy cơ tràn bộ nhớ (Out-of-Memory). Các kết quả thực nghiệm mô phỏng trên nền dữ liệu Non-IID cực đoan đã chứng minh sự vượt trội tuyệt đối của FedKDL cả về hiệu suất nhận diện mAP lẫn tốc độ hội tụ so với các hệ thống phân tán truyền thống.

\subsection{Hướng nghiên cứu tương lai (Future Works)}\label{sec1.7.2}

Mặc dù FedKDL đã thiết lập một nền tảng kiến trúc vững chắc cho tác vụ nhận diện mục tiêu dưới nước, hệ sinh thái IoUT vẫn mở ra những không gian nghiên cứu sâu rộng mang tính cách mạng. Trong tương lai, chúng tôi định hướng mở rộng giới hạn của hệ thống trên bốn trục đột phá chính:

Thứ nhất, tích hợp Học tự giám sát (Self-Supervised Learning - SSL) để hướng tới tự chủ hoàn toàn. Bằng cách sử dụng một mạng cơ sở (backbone) đã được tiền huấn luyện qua cơ chế học đối chiếu trên dữ liệu biển thô, kết hợp với luồng tinh chỉnh phần dư LoRA của FedKDL, hệ thống sẽ kiến tạo một đường ống (pipeline) học tập khép kín. Điều này cho phép bầy đàn AUV tự sinh nhãn giả và liên tục tiến hóa dưới đáy đại dương mà không phụ thuộc vào dữ liệu gán nhãn đắt đỏ từ con người.

Thứ hai, ứng dụng Khoảng cách Earth Mover (Earth Mover's Distance - EMD) trong phân cụm động. Các thuật toán phân cụm hiện tại chủ yếu dựa trên khoảng cách địa lý thường bỏ qua sự mất cân bằng về mặt ngữ nghĩa (semantic skewness) của dữ liệu. Bằng cách sử dụng EMD để định lượng trực tiếp sự sai khác phân phối nhãn giữa các AUV, mạng lưới có thể linh hoạt tái cấu trúc cụm (dynamic clustering) nhằm cân bằng phân phối dữ liệu cục bộ, từ đó giảm thiểu triệt để hiện tượng trôi dạt trọng số (Client Drift).

Thứ ba, điều phối tài nguyên bằng Học tăng cường (Reinforcement Learning - RL). Trong bối cảnh đại dương luôn biến động, các tác tử RL có thể được nhúng vào Gateway hoặc Relay để liên tục tối ưu hóa chính sách lựa chọn AUV tham gia huấn luyện. Thông qua việc học hỏi các hàm phần thưởng (reward functions) dựa trên năng lượng pin còn lại, chất lượng tín hiệu sóng âm, và độ quan trọng của gradient, RL sẽ thay thế hoàn toàn các cơ chế lập lịch tĩnh để tối đa hóa hiệu suất toàn cục.

Cuối cùng, việc phát triển cơ chế Thích nghi Rank Động (Adaptive LoRA Rank) song hành cùng Giao thức Tổng hợp An toàn (Secure Aggregation) sẽ cung cấp "lớp giáp" cuối cùng cho hệ thống. Khả năng tự co giãn kích thước ma trận nén theo tỷ số tín hiệu trên nhiễu (SNR) và khả năng tự động cô lập các trọng số độc hại (Byzantine faults) sẽ củng cố tính linh hoạt và độ tin cậy tuyệt đối cho một hệ thống IoUT quy mô lớn.


\backmatter
\bmhead{Acknowledgements}
This work was supported by Hanoi University of Science and Technology.

\section*{Declarations}
\begin{itemize}
\item \textbf{Conflict of interest:} The authors declare no competing interests.
\item \textbf{Data availability:} Simulation code and data are available upon reasonable request.
\end{itemize}

\bibliography{sn-bibliography}   % tên file .bib không cần đuôi

\end{document}
